\documentclass[11pt,a4paper]{article}

\usepackage[french]{babel}
\usepackage[utf8]{inputenc}
\usepackage[T1]{fontenc}
\usepackage{amsmath,amssymb,amsthm}
\usepackage{geometry}
\geometry{margin=2.6cm}
\usepackage{booktabs}
\usepackage{array}
\usepackage{longtable}
\usepackage{enumitem}
\usepackage{xcolor}
\usepackage{hyperref}
\hypersetup{colorlinks=true, linkcolor=blue!50!black, citecolor=blue!50!black, urlcolor=blue!40!black}
\usepackage{parskip}
\usepackage{titlesec}
\usepackage{tcolorbox}
\tcbuselibrary{skins,breakable}

\newtcolorbox{encadre}[1][]{colback=blue!4!white, colframe=blue!35!black, boxrule=0.5pt, arc=2pt, breakable, title={#1}, fonttitle=\bfseries, left=6pt, right=6pt, top=4pt, bottom=4pt}

\newtcolorbox{mesure}[1][]{colback=green!5!white, colframe=green!30!black, boxrule=0.5pt, arc=2pt, breakable, title={#1}, fonttitle=\bfseries, left=6pt, right=6pt, top=4pt, bottom=4pt}

\newcommand{\eng}[1]{\emph{#1}}

\title{\Large Comprendre la campagne d'expérimentation sur la fidélité\\des approximations de la matrice de Fisher\\[4pt]\large Un déroulé complet, sans raccourcis techniques}
\author{}
\date{Septembre 2026}

\begin{document}
\maketitle

\begin{abstract}
\noindent Ce document explique, en partant de zéro, ce que mesure la campagne d'expérimentation
décrite dans \texttt{docs/reports/plan\_exp\_draft.md} (et son brouillon original
\texttt{plan\_exp\_draft\_v0.md}). Le projet AdaFisher, et les quatre variantes construites autour
de lui dans ce dépôt (K-FAC, EKFAC, TKFAC, TEKFAC), reposent tous sur une idée simple mais jamais
vraiment vérifiée expérimentalement : au lieu de calculer exactement la « courbure » du problème
d'apprentissage — une quantité bien trop grosse pour être manipulée telle quelle sur un réseau de
neurones réaliste — on la remplace par une version très compressée, bon marché à calculer et à
inverser. La question posée par cette campagne est directe : \emph{à quel point cette version
compressée ressemble-t-elle à l'originale, et où précisément perd-elle le plus d'information ?}
Ce document reconstruit, sans supposer de connaissances préalables, tous les objets mathématiques
en jeu, explique pourquoi chaque expérience est nécessaire, et rapporte ce qui a déjà été mesuré.
\end{abstract}

\tableofcontents

% =====================================================================
\section{Introduction et mise en contexte}
% =====================================================================

\subsection{D'où vient ce projet}

Le dépôt dans lequel s'inscrit ce plan développe un optimiseur, \textbf{AdaFisher}, et cinq façons
différentes (« modes ») de construire l'ingrédient central dont il a besoin à chaque pas
d'entraînement : une estimation de la courbure du problème, qui sert à corriger la direction et la
taille du pas de descente de gradient. Les cinq modes (\texttt{diag}, \texttt{kfac}, \texttt{ekfac},
\texttt{tkfac}, \texttt{tekfac}) ont déjà été implémentés, testés unitairement et comparés entre eux
sur plusieurs bancs d'essai (un auto-encodeur MNIST, des CNN et des transformeurs sur CIFAR-10, etc.).
Ce travail d'implémentation — huit « lots » numérotés dans \texttt{docs/reports/plan.md} — a répondu
à une première question : \emph{les cinq modes sont-ils correctement implémentés, au sens où chacun
calcule bien ce que sa formule mathématique prescrit ?} La réponse est oui, vérifiée par des
centaines de tests.

Mais une question plus profonde reste ouverte, et c'est elle que ce nouveau plan attaque : \emph{ces
formules elles-mêmes sont-elles de bonnes approximations de la quantité qu'elles prétendent
approcher ?} Un mode peut être implémenté à la perfection tout en étant, dans l'absolu, une très
mauvaise idée si l'hypothèse mathématique sur laquelle il repose ne tient pas sur les réseaux réels.
C'est exactement le genre de question que l'article original d'AdaFisher n'a testée que très
partiellement : sa propre validation (annexe B.2 du papier) se limite à un petit réseau jouet, une
« vraie Fisher » obtenue par tirage aléatoire (Monte-Carlo) plutôt que calculée exactement, et une
seule mesure (une erreur moyenne entre deux diagonales de matrices). La campagne décrite ici reprend
ce travail depuis le début, en le rendant beaucoup plus rigoureux : calcul exact (et non approché)
des quantités de référence, huit types de couches différents, plusieurs mesures complémentaires, et
un « plancher de bruit » qui dit explicitement à partir de quand une différence mesurée cesse d'être
significative.

\subsection{Ce que ce document est, et n'est pas}

Ce document \textbf{n'introduit aucune idée nouvelle} : il traduit, sans jargon et avec les
justifications mathématiques complètes, le contenu de \texttt{plan\_exp\_draft.md} (la version
actuellement suivie par le code, notée « v1 » ci-dessous) et de son brouillon original en français,
\texttt{plan\_exp\_draft\_v0.md} (noté « v0 »), qui définit les questions et hypothèses de recherche
dans leur forme la plus explicite. Là où v1 a resserré le périmètre de v0 (moins de modèles, moins de
métriques dans un premier temps, un sous-ensemble d'hypothèses priorisé), ce document le signale
clairement plutôt que de mélanger les deux versions silencieusement.

Une distinction à garder en tête tout du long : cette campagne \textbf{ne modifie rien} au code de
l'optimiseur (\texttt{src/adafisher\_modes/}) ni aux bancs d'essai (\texttt{benchmarks/}). Elle
construit un paquet complètement séparé, \texttt{fisher\_ref/}, qui se contente de \emph{lire} les
modèles entraînés (leurs poids sauvegardés à différents instants de l'entraînement) et de calculer,
sur ces instantanés figés, des quantités de référence exactes auxquelles comparer chaque
approximation. C'est un instrument de mesure, pas une nouvelle méthode d'entraînement.

% =====================================================================
\section{Les notions de base, construites depuis zéro}
% =====================================================================

\subsection{Pourquoi s'intéresser à autre chose que le gradient}

Entraîner un réseau de neurones consiste à ajuster ses paramètres $\theta$ pour rendre une fonction
de perte $\ell(\theta)$ aussi petite que possible. La méthode la plus simple, la descente de
gradient, ne regarde que la pente locale de cette fonction : à chaque pas, on avance dans la
direction où la perte diminue le plus vite, avec un pas dont la taille (le taux d'apprentissage) est
la même dans toutes les directions.

Le problème est que la perte n'est presque jamais également « raide » dans toutes les directions de
l'espace des paramètres. Imaginons un paysage en forme de vallée très allongée et étroite : la
direction perpendiculaire à la vallée est très raide (un petit déplacement fait beaucoup remonter la
perte), la direction le long de la vallée est presque plate. Un pas de gradient bien réglé pour ne
pas diverger dans la direction raide sera ridiculement petit dans la direction plate, et
l'entraînement progresse alors très lentement. L'information qui manque au gradient seul, c'est la
\textbf{courbure} : comment la pente elle-même change quand on se déplace. C'est exactement ce que
mesure la matrice hessienne (ou une de ses cousines, voir plus bas) : une matrice carrée de taille
$P\times P$ ($P$ = nombre de paramètres), dont l'entrée $(i,j)$ dit à quel point la pente selon le
paramètre $i$ change quand on bouge le paramètre $j$.

Connaître cette matrice permettrait de calculer, en un seul pas bien orienté et bien dimensionné, ce
que le gradient seul mettrait des centaines de petits pas hésitants à approcher. C'est la promesse
des méthodes dites de \emph{second ordre} (par opposition au gradient, qui est de \emph{premier
ordre}), dont la méthode du gradient naturel — l'ancêtre conceptuel direct d'AdaFisher — fait partie.

\subsection{Le problème d'échelle, et pourquoi on approxime}

Le problème pratique est immédiat : pour un réseau de quelques millions de paramètres, la matrice de
courbure complète a des millions au carré d'entrées — bien trop pour tenir en mémoire, et de toute
façon impossible à inverser (l'opération dont on a besoin pour transformer le gradient en un « bon »
pas) en un temps raisonnable. Toutes les méthodes pratiques de second ordre — y compris AdaFisher et
ses quatre variantes — sont donc des \textbf{compromis} : elles ne gardent qu'une petite partie de
l'information contenue dans cette énorme matrice, sous une forme bon marché à stocker et à inverser.
La question de cette campagne est de mesurer, précisément, ce que chaque compromis perd.

\subsection{La matrice de Fisher : une courbure liée à la vraisemblance}

Il existe plusieurs façons raisonnables de définir « la courbure » d'un problème d'apprentissage. La
famille étudiée ici tourne autour de la \textbf{matrice d'information de Fisher}, notée $F$, qui a
un statut particulier pour les modèles entraînés par maximum de vraisemblance (ce qui couvre la
quasi-totalité des cas pratiques : classification par entropie croisée, régression par erreur
quadratique).

Concrètement, pour un réseau qui transforme une entrée $x_n$ en une sortie $z_n$ (les « logits » en
classification), puis produit une distribution de probabilité sur les labels possibles,
$r(y\mid z_n)$, la matrice de Fisher s'écrit
\[
F \;=\; \frac1N\sum_{n=1}^N J_n^\top \Lambda_n J_n,
\qquad J_n = \frac{\partial z_n}{\partial\theta}, \qquad
\Lambda_n = \mathrm{Cov}_{y\sim r(\cdot\mid z_n)}\!\big[\nabla_z \log r(y\mid z_n)\big].
\]
En mots : $J_n$ décrit comment la sortie du réseau bouge quand on bouge chaque paramètre (c'est une
information purement architecturale, calculable sans référence aux labels), et $\Lambda_n$ décrit à
quel point la distribution de probabilité que le modèle produit lui-même est sensible à un
changement de sa propre sortie. Pour une classification softmax à $C$ classes de probabilités
$p_n=(p_{n,1},\dots,p_{n,C})$, cette matrice $\Lambda_n$ a une forme fermée très simple,
$\Lambda_n=\mathrm{diag}(p_n)-p_np_n^\top$ : c'est la matrice de covariance d'un vecteur one-hot
tiré selon $p_n$.

Un fait mathématique important, que ce plan cite systématiquement, est que cette matrice de Fisher
coïncide exactement avec une autre courbure très utilisée en pratique, la matrice de
Gauss--Newton généralisée (\eng{GGN}), \textbf{à condition} que la fonction de perte utilisée soit
liée à $r$ par ce qu'on appelle un « lien canonique » (Martens, 2020, section~9) : c'est vrai pour
l'entropie croisée sur des logits, l'entropie croisée binaire sur des logits, et l'erreur
quadratique — c'est-à-dire pour tous les cas que ce dépôt utilise réellement. C'est ce qui permet au
plan de parler indifféremment de « Fisher vraie » et de « GGN » pour les modèles étudiés.

\subsection{Deux façons de faire la moyenne : vraie Fisher et Fisher empirique}

La formule ci-dessus fait apparaître une espérance implicite : $\Lambda_n$ moyenne sur
\emph{toutes les valeurs possibles} de $y$ que le modèle jugerait plausibles pour l'entrée $x_n$,
pondérées par la probabilité que le modèle leur attribue lui-même. C'est ce qui rend le calcul de
$F$ coûteux : en pratique, on ne peut approcher exactement cette moyenne sur $C$ classes qu'en faisant
$C$ passes de rétropropagation par exemple (une par classe possible), ou en utilisant une formule
fermée quand elle existe (le cas softmax ci-dessus).

Il existe une alternative beaucoup moins chère : au lieu de moyenner sur toutes les étiquettes
imaginables, on utilise \textbf{uniquement} l'étiquette réellement observée $y_n$ pour chaque
exemple. On obtient alors ce qu'on appelle la \textbf{Fisher empirique} :
\[
\hat E = \frac1N\sum_n \nabla_\theta \ell_n \, \nabla_\theta \ell_n^\top,
\qquad \ell_n = -\log r(y_n\mid z_n).
\]
Cette quantité ne coûte qu'\textbf{une seule} passe de rétropropagation par exemple (celle qu'on
fait de toute façon pour calculer le gradient d'entraînement), contre $C$ pour la vraie Fisher. C'est
la raison structurelle pour laquelle \emph{tous} les optimiseurs pratiques — AdaFisher inclus — n'ont
jamais accès, dans la boucle d'entraînement réelle, qu'à une version de type « Fisher empirique »,
jamais à la vraie Fisher.

Le problème est que $F$ et $\hat E$ ne sont \textbf{pas} la même quantité, et peuvent différer
énormément, même quand le modèle est déjà bien entraîné (Kunstner, Balles \& Hennig, 2019) : à
l'optimum, le gradient d'un exemple individuel ne s'annule pas forcément, alors que la vraie
distribution du modèle, elle, devient très proche de la vraie distribution des données. Mesurer
précisément cet écart, indépendamment de toute approximation supplémentaire, est le premier axe de
mesure de cette campagne (voir la question Q1, section~\ref{sec:design}).

\subsection{Résumé intuitif}

\begin{encadre}[À retenir]
\begin{itemize}[leftmargin=1.4em]
\item La courbure exacte d'un réseau de neurones est une matrice bien trop grosse pour être
manipulée directement : toute méthode pratique en garde une version compressée.
\item Pour les modèles entraînés par maximum de vraisemblance, la courbure naturelle est la matrice
de Fisher $F$, qui coïncide (sous une condition presque toujours vérifiée ici) avec la matrice de
Gauss--Newton.
\item $F$ ($\Lambda$ moyennée sur les étiquettes imaginées par le modèle lui-même) est chère à
calculer ; $\hat E$ (Fisher empirique, qui n'utilise que l'étiquette réelle) est bon marché mais
biaisée. Les optimiseurs pratiques n'ont jamais accès qu'à une version de type $\hat E$.
\item Cette campagne calcule, sur des instantanés figés de modèles réels, $F$ et $\hat E$ de manière
\textbf{exacte} (sans aucune structure imposée), pour servir de référence à laquelle comparer chaque
approximation utilisée en pratique.
\end{itemize}
\end{encadre}

% =====================================================================
\section{Le « zoo » des approximations : comment on compresse la courbure}
\label{sec:zoo}
% =====================================================================

Pour chaque couche $\ell$ du réseau (une couche linéaire, une convolution, une couche de
normalisation, \dots), on peut isoler le bloc de la matrice de courbure qui ne concerne que les
paramètres de cette couche. Notons $a_{n,t}$ l'entrée de la couche pour l'exemple $n$ à la position
partagée $t$ (un token dans une séquence, une position spatiale dans une image — pour une couche
sans partage de poids, il n'y a qu'une seule position), avec un $1$ ajouté pour représenter le biais,
et $g_{n,c,t}$ le gradient qui redescend jusqu'à la sortie de cette couche (avant sa non-linéarité),
pour la colonne $c$ associée à une des sources décrites plus bas. Le gradient de la couche pour cet
exemple et cette colonne est $\mathcal G_{n,c}=\sum_t g_{n,c,t}\,a_{n,t}^\top$, et le bloc exact de
courbure de la couche s'écrit
\[
B_\ell = \frac1N\sum_{n,c} \mathrm{vec}(\mathcal G_{n,c})\,\mathrm{vec}(\mathcal G_{n,c})^\top .
\]
Ce bloc $B_\ell$ est \emph{exact} : c'est le sous-bloc diagonal de $F$ (si $c$ parcourt les colonnes
de la source « vraie Fisher », voir section~\ref{sec:sources}) ou de $\hat E$ (source « empirique »)
qui concerne la couche $\ell$. Toutes les approximations ci-dessous en sont des simplifications.

\subsection{Le bloc-diagonal exact : ignorer les couches entre elles}

La première simplification, la plus douce, consiste à garder $B_\ell$ \textbf{exactement} pour
chaque couche, mais à supposer que deux couches différentes n'ont aucune corrélation de courbure
entre elles (les blocs hors diagonale, entre couches, sont mis à zéro). C'est un compromis
raisonnable si l'information qui circule entre couches est réellement faible — une hypothèse que
cette campagne teste directement (métrique M8 et hypothèse HF5, plus bas).

\subsection{K-FAC : factoriser en produit de Kronecker}

L'idée centrale de K-FAC (Martens \& Grosse, 2015) est de remarquer que, \emph{si} les entrées
$a_{n,t}$ de la couche et les gradients $g_{n,c,t}$ qui en ressortent étaient statistiquement
indépendants, le bloc $B_\ell$ se factoriserait exactement en un \textbf{produit de Kronecker}
$A\otimes G$ de deux matrices bien plus petites : une matrice $A$ qui ne décrit que les corrélations
entre les entrées de la couche, et une matrice $G$ qui ne décrit que les corrélations entre les
gradients de sortie. Concrètement,
\[
A = \frac1{NT}\sum_{n,t} a_{n,t}a_{n,t}^\top, \qquad G = \frac1N\sum_{n,c,t} g_{n,c,t}g_{n,c,t}^\top,
\qquad K_{\mathrm{KFAC}} = A\otimes G.
\]
L'intérêt est majeur : si la couche a $d_{\rm in}$ entrées et $d_{\rm out}$ sorties, le bloc exact a
$(d_{\rm in}d_{\rm out})^2$ entrées, alors que $A$ et $G$ n'en ont que $d_{\rm in}^2+d_{\rm out}^2$ —
et surtout, l'inverse d'un produit de Kronecker se calcule comme le produit de Kronecker des
inverses, $\big(A\otimes G\big)^{-1}=A^{-1}\otimes G^{-1}$, ce qui rend le préconditionnement
abordable à chaque pas d'entraînement.

Le prix de cette économie est double, et les deux effets sont mesurés séparément par cette
campagne :
\begin{itemize}
\item pour une couche à \textbf{poids partagés} (une convolution, ou une couche linéaire appliquée à
chaque position d'une séquence comme dans l'attention), il faut d'abord décider comment traiter les
différentes positions $t$ : soit les traiter comme des exemples indépendants qu'on empile
(\eng{expand}, exact pour un réseau linéaire profond selon Eschenhagen et al., 2023), soit résumer
d'abord chaque exemple par une somme sur ses positions avant de corréler (\eng{reduce}, exact quand
l'agrégation finale est une moyenne pondérée, par exemple un \eng{pooling} moyen). Le choix
expand/reduce n'est pas neutre, et son erreur est mesurée par la décomposition « partage de poids »
de la section~\ref{sec:design} ;
\item même en ignorant le partage de poids, l'hypothèse d'\textbf{indépendance} entre entrées et
gradients de sortie n'est vraie qu'en général de façon approximative. L'écart entre le bloc exact
$B_\ell$ et sa version « produit de Kronecker le plus proche » est l'\textbf{erreur d'indépendance},
mesurée directement (métrique M7).
\end{itemize}

\subsection{EKFAC : corriger K-FAC dans sa propre base}

K-FAC choisit implicitement une direction privilégiée pour chaque paramètre de la couche : les
vecteurs propres du produit de Kronecker $A\otimes G$, c'est-à-dire les produits croisés des
vecteurs propres de $A$ (notés $Q_A$) et de $G$ (notés $Q_B$). Le raffinement EKFAC (George et al.,
2018) garde exactement cette même base de directions, mais réajuste la « force » attribuée à chacune
de ces directions de façon \textbf{optimale}, en la recalculant directement à partir des données :
\[
s_{ij} = \frac1N\sum_{n,c}\Big[(Q_G^\top \mathcal G_{n,c} Q_A)_{ij}\Big]^2 .
\]
Cette correction est peu coûteuse (elle ne change pas la base, seulement des coefficients), et elle
garantit par construction que la trace totale du bloc approché reste égale à celle du bloc exact
$B_\ell$ — une propriété que K-FAC lui-même n'a pas. Intuitivement, EKFAC dit : « je garde les mêmes
directions que K-FAC, parce que les calculer autrement serait cher, mais je ne suppose plus que leur
importance suit exactement le produit des deux facteurs séparés — je la mesure directement ».

\subsection{TKFAC et TEKFAC : une autre factorisation, qui préserve l'énergie totale}

TKFAC (Gao et al., 2021) part d'une factorisation différente, construite pour garantir exactement
une propriété précise : la trace du bloc de courbure approché (une mesure de son « énergie » totale,
la somme de ses valeurs propres) est \textbf{égale}, par construction, à la trace du bloc exact
$B_\ell$ — et ce sans avoir besoin de recalculer quoi que ce soit après coup, contrairement à EKFAC.
TEKFAC applique à TKFAC le même raffinement qu'EKFAC applique à K-FAC : garder la base propre issue
de TKFAC, mais réajuster optimalement l'importance de chaque direction.

\subsection{La diagonale : le compromis le plus radical}

La version la plus économique de toutes consiste à ne garder que la diagonale du bloc de courbure,
c'est-à-dire à supposer qu'il n'existe \textbf{aucune} corrélation entre deux paramètres différents,
même à l'intérieur d'une même couche — chaque paramètre est traité indépendamment de tous les
autres. C'est ce que fait le mode \texttt{diag} d'AdaFisher, avec une propriété remarquable : la
diagonale d'un produit de Kronecker $A\otimes G$ est exactement le produit de Kronecker des
diagonales, $\mathrm{diag}(A)\otimes\mathrm{diag}(G)$. Autrement dit, l'estimateur « brut »
d'AdaFisher \textbf{est} la diagonale de K-FAC. Il existe donc deux chemins différents pour
l'obtenir : le calculer directement comme diagonale du bloc exact $B_\ell$, ou le calculer comme
diagonale du produit de Kronecker $A\otimes G$. Ces deux chemins ne coïncident pas exactement, et leur
écart, entrée par entrée, vaut précisément
\[
\mathbb E[a_j^2 g_i^2] - \mathbb E[a_j^2]\,\mathbb E[g_i^2] \;=\; \mathrm{Cov}(a_j^2, g_i^2),
\]
c'est-à-dire le \textbf{biais d'indépendance restreint à la diagonale seule} — une quantité que la
campagne mesure directement, car elle est distincte de la perte, bien plus large, de toute
l'information hors diagonale.

Ce qu'AdaFisher utilise réellement dans l'optimiseur n'est pas cette diagonale « brute », mais une
version encore transformée pour la rendre numériquement stable et utilisable pas après pas :
normalisée entre 0 et 1 (« min-max »), lissée dans le temps par une moyenne mobile exponentielle
(EMA) qui combine l'estimation du pas courant avec la mémoire des pas précédents, et enfin amortie
par un petit terme $\lambda$ pour éviter de diviser par une valeur trop proche de zéro. Cette version
pleinement opérationnelle est appelée « AF-op » dans le plan, par opposition à la diagonale « brute »,
« AF-raw ». La différence entre les deux est précisément l'objet du deuxième protocole de mesure,
P2 (section~\ref{sec:protocoles}).

\subsection{Le témoin diagonal libre : Adam}

Un point de comparaison utile, gratuit, est le second moment que l'optimiseur Adam maintient déjà
pour chaque paramètre (la moyenne mobile des carrés du gradient). Ce n'est pas conçu comme une
approximation de courbure, mais des travaux récents (Li, Dangel, Tam \& Raffel, 2025) soutiennent
que toute méthode diagonale « sophistiquée » devrait au moins faire mieux que ce témoin gratuit, sans
quoi la sophistication ne sert à rien. Les campagnes de bancs d'essai de ce dépôt font tourner Adam
et AdamW comme deux des sept « bras » comparés systématiquement (voir section~\ref{sec:checkpoints}),
ce qui fournit ce témoin sans coût supplémentaire.

\subsection{Un cas particulier important : les couches de normalisation}

Les couches de normalisation (\texttt{LayerNorm}, \texttt{GroupNorm}, \texttt{BatchNorm2d} en mode
évaluation) ont deux paramètres par canal, une échelle $\gamma$ et un décalage $\beta$, appliqués à
une entrée déjà normalisée $\hat x_{n,t}$ (de moyenne nulle et de variance unité). Leurs gradients
exacts par exemple sont simples : $\nabla_\gamma=\sum_t g_t\odot\hat x_t$ et
$\nabla_\beta=\sum_t g_t$.

Un résultat mathématique établi pour AdaFisher (la proposition~3.1 de son article) montre que le
bloc de courbure exact d'une telle couche \textbf{n'a pas} la structure « produit de Kronecker »
des couches précédentes, mais une structure différente, dite de Hadamard (un produit terme à
terme entre deux matrices, plutôt qu'un produit de Kronecker). Faire entrer cette structure
différente dans le même moule que les quatre approximations précédentes (qui attendent toutes un
produit de Kronecker $A\otimes G$) demande un choix de conception explicite — un petit facteur
scalaire choisi pour être la meilleure approximation possible (au sens des moindres carrés,
« optimale au sens de Frobenius ») de la vraie structure de Hadamard. Ce choix, documenté et
délibéré, est un des points que la campagne vérifie chiffre en main (c'était l'hypothèse HF4 de v0,
voir plus bas), et le repère a en réalité déjà été largement tranché analytiquement au moment
d'implémenter les modes (lot 5) : ce que cette campagne ajoute est la mesure \emph{empirique}, sur
un bloc exact $2C\times 2C$ toujours calculable (même dans le régime le plus contraint en mémoire),
de l'écart entre trois variantes : le bloc exact complet (avec le couplage entre $\gamma$ et
$\beta$), le bloc exact « séparé » (ce que fait AdaFisher, qui jette ce couplage), et la structure de
Hadamard approchée réellement implémentée.

\subsection{Les autres cas particuliers : \eng{embeddings}, \eng{LoRA}, la tête}

\begin{itemize}
\item \textbf{Position/embedding} (\texttt{pos\_embed}, \texttt{cls\_token}) : ce sont des paramètres
« nus », qui n'appartiennent à aucune couche standard reconnue par les hooks de l'optimiseur — leur
courbure y est donc traitée comme l'identité (aucune information de courbure n'est utilisée pour
eux), un choix qu'il faut savoir expliciter plutôt que de le découvrir accidentellement.
\item \textbf{LoRA} (adaptateurs de \eng{fine-tuning} léger, non encore couverts par les modèles
actuellement construits) : un adaptateur s'écrit $W_0+BA$, un produit de deux petites matrices. La
courbure présente naturellement un couplage entre les deux adaptateurs $A$ et $B$ (le produit $BA$
est bilinéaire), que la métrique M8 (couplage inter-blocs) est destinée à quantifier si ce cas est
ajouté à la campagne.
\item \textbf{La tête de classification} : c'est une couche linéaire ordinaire, sans partage de
poids, mais avec une particularité : la matrice $\Lambda(p)$ qui pondère la sortie dépend
elle-même de l'entrée $a$ de cette couche (par l'intermédiaire de $p=\mathrm{softmax}(z)$).
L'hypothèse d'indépendance entre entrée et gradient, sur laquelle repose K-FAC, y est donc fausse
« par construction » — c'est un point de mesure informatif plutôt qu'un problème à corriger.
\end{itemize}

% =====================================================================
\section{Les trois sources du signal rétropropagé}
\label{sec:sources}
% =====================================================================

Dans la formule de $B_\ell$ ci-dessus, l'indice $c$ parcourt un ensemble de « colonnes » qui dépend
de la source $\sigma$ choisie pour le vecteur $g$ — c'est-à-dire de la façon dont on choisit de faire
la moyenne sur les étiquettes possibles évoquée en section~2.4 :

\begin{itemize}
\item \textbf{type-2} (la vraie Fisher) : $c=1,\dots,C$, une colonne par classe possible, en
utilisant la racine fermée du softmax $S_n[:,c]=\sqrt{p_c}\,(e_c-p_n)$, qui vérifie exactement
$S_nS_n^\top = \Lambda_n$. C'est coûteux ($C$ rétropropagations par exemple) mais exact, sans aucun
tirage aléatoire.
\item \textbf{Monte-Carlo} ($\mathrm{MC}_K$) : au lieu de considérer toutes les classes, on en tire
$K$ au hasard selon la distribution $p_n$ du modèle lui-même, et on utilise le vecteur
$K^{-1/2}(p_n - e_{\tilde y})$ pour chaque tirage $\tilde y$. C'est une estimation \textbf{stochastique}
de la vraie Fisher (elle y converge quand $K$ augmente, à la vitesse $O(K^{-1/2})$ classique de tout
estimateur Monte-Carlo), moins chère quand $K\ll C$. Le facteur $K^{-1/2}$ est indispensable : sans
lui, l'estimateur ne converge tout simplement pas vers la bonne quantité.
\item \textbf{empirique} : $c=1$, une seule colonne, en utilisant directement l'étiquette réellement
observée, $p_n-e_{y_n}$. C'est ce que tout optimiseur pratique utilise réellement.
\end{itemize}

Chaque structure d'approximation (K-FAC, EKFAC, \dots) peut être appliquée à chacune de ces trois
sources : \textbf{« structure » et « source » sont deux axes de variation indépendants}, et c'est
précisément ce croisement qui constitue le cœur du design expérimental (section suivante).

% =====================================================================
\section{Le design expérimental : croiser structure, source et référence}
\label{sec:design}
% =====================================================================

\subsection{Le tableau à trois entrées}

Pour chaque couche et chaque instant de l'entraînement, on calcule un écart $d\big(S(\sigma), R\big)$
entre :
\begin{itemize}
\item une \textbf{structure} $S$ : bloc complet, bloc-diagonal exact, K-FAC-expand, K-FAC-reduce,
EKFAC, TKFAC, diagonale-Kronecker d'AdaFisher, diagonale exacte, (et pour les normalisations :
Hadamard ou la version réellement implémentée) ;
\item une \textbf{source} $\sigma$ : type-2, Monte-Carlo à $K\in\{1,4,16\}$ tirages, ou empirique ;
\item une \textbf{référence} $R$ : la vraie Fisher $F$ ou la Fisher empirique exacte $\hat E$.
\end{itemize}

Certaines cellules de ce tableau à trois entrées répondent directement, chacune, à une question
précise :

\begin{center}
\begin{tabular}{@{}p{4.6cm}p{5.2cm}p{2.6cm}@{}}
\toprule
Cellule mesurée & Ce qu'elle isole & Question \\
\midrule
$d(\hat E, F)$ & l'erreur de \textbf{source} seule (aucune structure imposée) & Q1 \\
$d(S(\text{type-2}), F)$ & l'erreur de \textbf{structure} seule, contre la vraie Fisher & Q2 \\
$d(S(\text{emp}), \hat E)$ & l'erreur de \textbf{structure} seule, contre la Fisher empirique & Q2 \\
$d(S(\text{emp}), F)$ & le chemin \textbf{réellement emprunté} par AdaFisher en pratique & Q3 \\
$d(S(\mathrm{MC}_K), F)$ & structure + bruit d'échantillonnage Monte-Carlo & HF6 \\
tout ce qui précède, par type de couche & tout, désagrégé & Q4 \\
\bottomrule
\end{tabular}
\end{center}

En clair : \textbf{Q1} demande « si je ne fais \emph{aucune} approximation de structure, combien
est-ce que je perds simplement en remplaçant les vraies étiquettes imaginées par le modèle par les
étiquettes réelles observées ? ». \textbf{Q2} demande l'inverse : « à source fixée, combien
est-ce que je perds en compressant la structure (diagonale, Kronecker, \dots) ? ». \textbf{Q3}
combine les deux dans l'ordre où AdaFisher les subit réellement en pratique (structure compressée
\emph{et} source empirique), et compare le résultat à la vraie Fisher. \textbf{Q4} demande si les
réponses aux trois questions précédentes dépendent du type de couche (une couche linéaire simple
n'est pas forcément dans le même cas qu'une tête d'attention ou qu'une normalisation).

\subsection{L'erreur totale se décompose-t-elle proprement ? (Q3)}

Une question naturelle, une fois Q1 et Q2 mesurées séparément, est de savoir si l'erreur totale du
chemin réel (Q3) est simplement leur somme. Pour une vraie distance mathématique (comme l'erreur de
Frobenius ou l'erreur spectrale, voir section suivante), l'inégalité triangulaire garantit
\emph{seulement} une borne supérieure :
\[
\underbrace{d(S(\text{emp}), F)}_{\text{erreur totale}} \;\le\;
\underbrace{d(\hat E, F)}_{\text{erreur de source}} + \underbrace{d(S(\text{emp}), \hat E)}_{\text{erreur de structure}}.
\]
Ce que cette campagne rapporte, c'est \textbf{l'écart entre les deux membres} de cette inégalité :
si l'erreur totale mesurée est très inférieure à la somme des deux erreurs séparées, cela signifie
que les deux sources d'erreur ont tendance à \emph{se compenser} l'une l'autre (une bonne nouvelle
accidentelle) ; si elle est proche de la somme, les deux erreurs \emph{s'additionnent} sans se gêner
ni s'annuler. Pour la métrique de divergence de Kullback--Leibler (section~\ref{sec:metriques},
métrique M3), qui n'est pas une distance mathématique au sens strict (elle n'est pas symétrique), on
se contente de rapporter les trois termes côte à côte, sans chercher à les additionner.

\subsection{La décomposition propre aux couches à poids partagés}
\label{sec:partage}

Pour une couche à poids partagés (une convolution appliquée à chaque position spatiale, une couche
linéaire appliquée à chaque position d'une séquence dans l'attention), le passage du bloc exact
$B_\ell$ à K-FAC-expand se fait en deux étapes distinctes, dont chacune a un coût mesurable
séparément :
\[
B_\ell \;\xrightarrow[\text{croisés }t\ne t'\text{ ignorés}]{\text{erreur de \textbf{partage de poids}}}\;
B_\ell^{\rm expand} \;\xrightarrow[\ ]{\text{erreur d'\textbf{indépendance}}}\; \text{K-FAC-expand}.
\]
La première flèche consiste à ignorer les corrélations entre deux positions différentes de la même
couche (par exemple, entre deux pixels différents d'une même image traversant la même convolution) ;
la seconde est l'hypothèse d'indépendance entre entrées et gradients de sortie déjà décrite pour
K-FAC en général. Les mesurer séparément permet de savoir laquelle des deux hypothèses coûte le
plus cher sur une architecture donnée — une information directement exploitable pour décider où
investir un effort de modélisation plus fin. C'est maintenant fait, et la réponse est nette :
sur les trois modèles à poids partagés mesurés, la première flèche emporte à elle seule une
fraction médiane de $0{,}971$ du bloc exact, la seconde $0{,}45$ à $0{,}60$ de ce qui reste
(section~\ref{sec:lot3}).

% =====================================================================
\section{Les trois régimes de calcul, dictés par la mémoire disponible}
\label{sec:regimes}
% =====================================================================

Calculer $F$ ou $\hat E$ de façon exacte, sans aucune structure imposée, revient à former une
matrice $P\times P$ pleine — le même problème d'échelle que celui qui a motivé, en premier lieu,
toutes les approximations décrites ci-dessus. La campagne le résout en distinguant trois régimes de
calcul selon la taille du réseau étudié.

\subsection{L'astuce commune : voir $F$ comme un produit de deux petites choses}

Le point de départ commun est une écriture différente de $F$. On empile, pour chaque exemple sonde
$n$ et chaque colonne $c$ de la racine $\Lambda_n^{1/2}$ (par exemple les $C$ colonnes de la racine
softmax fermée $S_n$ décrite plus haut), une ligne $N^{-1/2}(\Lambda_n^{1/2}e_c)^\top J_n$ dans une
grande matrice $U\in\mathbb R^{m\times P}$, avec $m=N\cdot(\text{rang de la racine})$. On a alors
exactement
\[
F = U^\top U.
\]
Chaque ligne de $U$ ne coûte qu'\emph{une seule} rétropropagation d'un vecteur de sortie. C'est cette
écriture qui permet de choisir, selon $P$ et $m$, la stratégie de calcul la moins coûteuse.

\subsection{Régime A — tout former explicitement}

Quand le réseau est assez petit ($P$ de l'ordre de quelques dizaines de milliers de paramètres, la
limite exacte étant $P_{\max}\approx\sqrt{B/24}$ où $B$ est la mémoire GPU disponible en octets —
soit environ 28\,900 paramètres sur 20 Go, ou seulement 20\,400 sur les 10 Go du plus petit
découpage matériel utilisé pour les entraînements), on peut se permettre de calculer $U^\top U$ pour
de vrai et d'obtenir une matrice $P\times P$ que l'on stocke intégralement en mémoire, en double
précision (\eng{fp64}) — une précision plus élevée que celle utilisée pendant l'entraînement, parce
que le spectre des valeurs propres de $F$ s'étend sur plus de huit ordres de grandeur, et la
précision simple habituelle (\eng{fp32}) ne suffit pas à représenter fidèlement les valeurs propres
les plus petites. C'est le seul régime où l'on peut calculer \textbf{tout} exactement : le spectre
complet, les blocs entre couches différentes, toutes les métriques sans aucune approximation
supplémentaire. Le prix est le nombre limité de modèles qui rentrent dans cette catégorie, et le fait
qu'il faut aussi choisir un nombre d'exemples sondes $N$ suffisamment grand pour que $F$ puisse être
de rang plein (sinon toute métrique qui passe par un inverse serait dominée par le petit terme
d'amortissement ajouté pour rendre cet inverse possible, plutôt que par le contenu réel de $F$).

\subsection{Régime B — ne jamais former la grande matrice}

Pour des modèles plus grands (des centaines de milliers de paramètres), former $U$ dans son
intégralité serait déjà trop coûteux en mémoire. Le régime B ne forme \textbf{jamais} $U$ ni $F$
entiers : il garde, couche par couche, un « Gram » beaucoup plus petit,
$K_\ell = U_\ell U_\ell^\top \in\mathbb R^{m\times m}$ (une matrice de taille liée au nombre
d'exemples sondes plutôt qu'au nombre de paramètres de la couche), et les statistiques brutes de
couche (les $a_{n,t}$ et $g_{n,c,t}$ eux-mêmes). Une identité algébrique classique (l'identité de
Woodbury, la même qui rend les filtres de Kalman efficaces) permet alors de calculer, sans jamais
former de matrice $P\times P$, la plupart des métriques d'intérêt : normes de Frobenius, traces,
inverses amortis, divergences de Kullback--Leibler. Deux façons pratiques de construire ce Gram par
couche existent, à choisir couche par couche selon celle qui coûte le moins cher : la matérialiser
directement, ou passer par une astuce dite « \eng{ghost} » qui l'obtient sans jamais former de
gradient individuel par exemple, empruntée aux méthodes de confidentialité différentielle.

Le prix de ce régime est différent de celui du régime A : $F$, construite à partir d'un nombre fini
$m$ d'exemples sondes, a un \textbf{rang} au plus $m$, qui peut être bien plus petit que $P$. Les
métriques qui ne passent pas par un inverse (Frobenius, direction) restent exactes ; celles qui en
passent par un mélangent deux effets qu'il faut soigneusement démêler : l'erreur réelle dans la
partie de l'espace que $F$ « voit » réellement, et une courbure \textbf{inventée} par une
approximation de rang plein comme K-FAC là où $F$ n'a strictement aucune information (son noyau).
C'est pourquoi toute conclusion tirée du régime B est systématiquement présentée en fonction d'un
paramètre d'amortissement $\lambda$ balayé sur plusieurs ordres de grandeur, et confrontée, quand
c'est possible, au régime A qui n'a pas ce problème.

\subsection{Régime C — sans jamais former de matrice du tout}

Pour les très grands modèles (un ResNet-50 complet, par exemple), même les Grams par couche du
régime B deviendraient trop gros. Le régime C ne calcule \textbf{que des produits} de la matrice de
courbure par un vecteur ($Fv$), obtenus en une passe avant plus une passe arrière, sans jamais
matérialiser aucune matrice intermédiaire. Cela donne accès à des quantités comme la norme spectrale
ou la trace, calculées par des algorithmes itératifs (Lanczos, gradient conjugué), au prix d'une
petite erreur stochastique ou d'approximation contrôlée. Ce régime, bien que prévu dans le plan, est
optionnel et repoussé aux derniers lots d'implémentation : il apporte un modèle de plus au prix
d'un jeu de métriques plus restreint et de moins bonne précision.

\subsection{Décision prise pour cette campagne}

Un outil externe existant, \texttt{curvlinops}, aurait pu fournir directement les opérateurs
« sans matrice » du régime C et servir d'oracle indépendant pour valider le reste du code. La
décision, prise et documentée au premier lot d'implémentation, est de \textbf{ne pas l'adopter}
(il n'est disponible ni dans l'environnement du projet ni sur le cluster de calcul utilisé) : à la
place, un oracle indépendant plus simple, construit directement à partir des outils de
différentiation automatique de PyTorch, a été écrit spécifiquement pour vérifier le code de la
campagne, et il concorde avec l'implémentation principale à $3{,}8\times 10^{-16}$ près sur des
vecteurs aléatoires — une précision largement suffisante pour servir de garde-fou.

% =====================================================================
\section{Les modèles étudiés, et le rôle de chacun}
% =====================================================================

Le tableau suivant reprend les huit dossiers de bancs d'essai (déjà construits et testés pour les
huit lots d'implémentation de l'optimiseur) réellement utilisés par cette campagne, avec le nombre
exact de paramètres de chacun et le régime de calcul dans lequel il tombe.

\begin{center}
\small
\begin{tabular}{@{}p{0.8cm}p{2.6cm}p{1.4cm}p{1.1cm}p{4.3cm}p{4.5cm}@{}}
\toprule
Id & Modèle & $P$ & Régime & Types de couches & Rôle dans la campagne \\
\midrule
A1 & MLP + LayerNorm (MNIST) & 26\,634 & A & Linear sans partage, LayerNorm, tête & le cas témoin, le plus propre pour la théorie de K-FAC \\
A2 & CNN + GroupNorm/BatchNorm (CIFAR-10) & 24\,458 & A & Convolution (partage spatial), normalisation, tête & reproduit le banc jouet de l'article AdaFisher, mais contre une Fisher exacte \\
A3 & ViT miniature (CIFAR-10) & 21\,098 & A & \eng{patch-embed}, attention QKV fusionnée, MLP, LayerNorm, tête & un cadre « \eng{reduce} » propre (agrégation par moyenne) \\
B0 & Auto-encodeur MNIST & 2\,837\,314 & B & 8 couches linéaires, sigmoïdes & le banc \textbf{principal} du projet, celui où les cinq modes se bloquent (section~\ref{sec:stall}) \\
B1 & CCT-2/3$\times$2 (CIFAR-10) & 283\,723 & B & tokenizer convolutionnel, transformeur, LayerNorm, \eng{pooling} de séquence & le modèle propre à l'article AdaFisher \\
B2 & ResNet-20 (CIFAR-10) & 269\,722 & B & convolution, BatchNorm & le petit ResNet de la famille testée par AdaFisher \\
B4 & ViT à largeur réaliste (CIFAR-10) & 2\,693\,578 & B, en flux & transformeur $d=192$, 6 blocs & vérifie que les conclusions tirées sur A3/B1 tiennent à une largeur réaliste \\
C & ResNet-50 (CIFAR-10) & 23\,520\,842 & C uniquement & convolutions profondes, BatchNorm & optionnel, régime sans matrice seulement \\
\bottomrule
\end{tabular}
\end{center}

Deux modèles envisagés dans le brouillon initial (un modèle de langage miniature de type GPT, et un
modèle de \eng{fine-tuning} léger par LoRA sur RoBERTa) ne sont \textbf{délibérément pas construits}
pour l'instant : ce sont les deux seuls cas qui porteraient des types de couches encore absents du
reste de la campagne (un \eng{embedding} de tokens, une attention causale, des adaptateurs LoRA), et
ils restent les premiers candidats pour un futur dossier de modèle, une fois que la carte principale
sera stabilisée sur les huit modèles ci-dessus.

Trois précisions structurelles à garder en tête, héritées directement de la façon dont ces modèles
ont été construits pour les besoins de l'optimiseur, avant même cette campagne :
\begin{itemize}
\item Dans A2, la variante \texttt{GroupNorm} n'est \textbf{pas} un type de couche reconnu par
l'optimiseur AdaFisher : ses paramètres reçoivent, dans l'optimiseur, un traitement équivalent à
« pas de courbure du tout » (préconditionneur identité). Ce n'est pas un oubli mais un choix
délibéré du banc, qui donne à la campagne un cas de contrôle gratuit : une normalisation dont le
bloc exact est parfaitement calculable, mais dont la contrepartie « opérationnelle » côté optimiseur
est simplement absente.
\item Dans A3, B1 et B4, l'attention expose une seule couche linéaire fusionnée qui calcule en un
bloc les trois projections \eng{query}, \eng{key} et \eng{value} — un choix d'implémentation
(hérité de la bibliothèque utilisée) qui empêche, pour l'instant, de mesurer séparément l'erreur de
K-FAC sur chacune des trois projections individuellement ; seule la version « bloc fusionné, entrée
partagée » est mesurable avec les modèles actuels.
\item Les vecteurs de position (\texttt{pos\_embed}) et le jeton de classification
(\texttt{cls\_token}), là où ils existent, sont des paramètres « nus » n'appartenant à aucune couche
reconnue : ils sont bien mis à jour pendant l'entraînement, mais leur courbure y est traitée comme
l'identité — un fait que le registre de classification des paramètres de la campagne doit
explicitement reconnaître plutôt que de le laisser passer inaperçu.
\end{itemize}

% =====================================================================
\section{Les deux protocoles de mesure : ce qui est possible contre ce qui se passe réellement}
\label{sec:protocoles}
% =====================================================================

Une même approximation (disons K-FAC) peut être évaluée de deux façons très différentes, et cette
campagne les distingue systématiquement.

\subsection{Protocole P1, structurel : la meilleure version possible}

Dans le protocole P1, on prend un jeu de poids $\theta$ figé (un instantané de l'entraînement), un
ensemble fixe d'exemples sondes, et on calcule chaque approximation \textbf{sans} aucune des
approximations supplémentaires que l'optimiseur ajoute en pratique pour tenir sur la durée d'un
entraînement complet : pas de lissage temporel (pas de moyenne mobile), pas de renormalisation
min-max, et un amortissement $\lambda$ balayé systématiquement plutôt que fixé à une seule valeur.
P1 répond à la question : \emph{« dans le meilleur des cas, quelle est la fidélité maximale que
cette famille d'approximation peut atteindre ? »}

\subsection{Protocole P2, opérationnel : ce que l'optimiseur utilise vraiment}

Dans le protocole P2, on va lire, directement dans l'état interne de l'optimiseur réel en train de
s'exécuter, la version de la courbure telle qu'elle sert \textbf{réellement} à corriger chaque pas :
après le lissage temporel (moyenne mobile exponentielle, avec un poids de 0,92 donné à l'historique
et 0,08 au facteur du mini-batch le plus récent, à chaque mise à jour tous les 100 pas), après la
renormalisation min-max (pour le mode \texttt{diag}), avec l'amortissement réellement utilisé en
entraînement ($\lambda=10^{-3}$), avec les couches de normalisation en mode entraînement (donc avec
un comportement légèrement différent du mode évaluation utilisé pour toutes les références exactes),
et avec le \eng{dropout} actif. On compare ensuite cet état réel aux références exactes, calculées
au \textbf{même} $\theta$. P2 répond à la question : \emph{« que se passe-t-il réellement, une fois
tous les compromis pratiques empilés les uns sur les autres ? »}

La comparaison entre P1 et P2 est directement ce qui permet de répondre à l'hypothèse HF7 (plus
bas) : si P2 est nettement pire que P1, la faiblesse ne vient pas de la structure mathématique
choisie (Kronecker, diagonale, \dots) mais des compromis d'ingénierie ajoutés autour d'elle (le
lissage temporel, la normalisation, l'amortissement) — une conclusion qui oriente la recherche vers
des pistes très différentes (par exemple, un lissage temporel plus fidèle à la géométrie du
problème, plutôt qu'une nouvelle structure d'approximation).

\subsection{Le problème du « réchauffage » (\eng{re-warm})}

Un obstacle pratique complique P2 : les instantanés sauvegardés pendant l'entraînement ne contiennent
que les \textbf{poids} du réseau, pas la mémoire interne de l'optimiseur (les moyennes mobiles
accumulées). Pour mesurer P2 à un instant donné de l'entraînement, il faut donc repartir des poids
sauvegardés, reconstruire un optimiseur neuf, et le laisser tourner un certain nombre de pas
supplémentaires pour que sa mémoire interne se « réchauffe » et redevienne représentative de ce
qu'elle aurait été si l'entraînement avait continué sans interruption — avant de faire la mesure.

Combien de pas faut-il pour ce réchauffage ? La règle de mise à jour de la moyenne mobile multiplie
l'ancien état par $0{,}92$ et y ajoute $0{,}08$ fois la nouvelle observation, à chaque intervalle de
100 pas (\texttt{TCov}). Après $k$ intervalles, la contribution de \emph{tout ce qui existait avant
le réchauffage} est multipliée par $0{,}92^k$ — ou, lu autrement, la composition des poids successifs
donne $0{,}92,\ 0{,}92\times0{,}08,\ \dots$, c'est-à-dire que \textbf{92\,\% de l'état opérationnel
n'est jamais que le facteur du tout dernier mini-batch}. Un point de correction important, découvert
en creusant cette arithmétique : la longueur de réchauffage nécessaire n'est pas fixée par la
condition naïve « $0{,}08^k$ doit être négligeable devant 1 », mais par la condition, plus stricte,
« $0{,}08^k$ doit être négligeable devant $\lambda$ » — parce que chaque optimiseur initialise sa
mémoire avec la matrice identité au tout premier pas, et qu'un réchauffage trop court laisse
subsister une trace de cette identité initiale qui agit comme un \textbf{amortissement parasite},
s'ajoutant au $\lambda$ voulu. À $k=3$ intervalles (300 pas), cette trace parasite
($0{,}08^3\approx 5\cdot10^{-4}$) vaut déjà la moitié de $\lambda=10^{-3}$, faussant le
préconditionneur appliqué de 12 à 87\,\% selon le mode. Il faut au moins $k=10$ intervalles (1000
pas) pour que cette contamination retombe à un niveau négligeable.

\begin{mesure}[Résultat mesuré — fidélité du réchauffage, \texttt{cnn\_gn\_cifar}, 11 septembre 2026]
Une expérience dédiée a comparé, à partir d'un même point $\theta^\star$ issu d'un vrai
entraînement de 2\,000 pas, trois optimiseurs neufs réchauffés sur des ordres de mini-batchs
différents : deux avec les poids figés (leur écart mutuel donne le \textbf{plancher de bruit}
propre à l'estimateur, dû uniquement au tirage aléatoire des mini-batchs), un avec les poids qui
continuent d'évoluer (donnant le terme de « désuétude », c'est-à-dire l'effet du déplacement de
$\theta$ pendant le réchauffage). Trois conclusions, mesurées et non supposées :
\begin{enumerate}[leftmargin=1.3em]
\item Après un réchauffage suffisant (au moins 10 intervalles, soit 1000 pas), chaque mode se
retrouve à seulement 1,3 à 4,5 fois son propre plancher de bruit — c'est-à-dire pratiquement
indiscernable d'une deuxième réalisation aléatoire de la même expérience. C'est ce qui rend inutile
de relancer toute la campagne d'entraînement uniquement pour conserver l'état de l'optimiseur : le
réchauffage suffit.
\item Le terme de désuétude (l'effet du déplacement des poids pendant le réchauffage) est,
lui, invisible : il est du même ordre que le bruit lui-même, ce qui n'a rien de surprenant puisque
la moyenne mobile ne donne de toute façon jamais plus de 8\,\% de poids à l'historique lointain.
\item Les facteurs bruts eux-mêmes restent, quelle que soit la longueur de réchauffage, à 2 à 4
fois leur propre plancher de bruit : ce ne sont, par nature, que des estimateurs à forte variance,
et relancer un entraînement complet ne ferait que tirer une nouvelle réalisation de cette même
variance, pas une valeur « plus vraie ».
\end{enumerate}
\end{mesure}

Un interrupteur optionnel, désactivé par défaut, permet désormais de sauvegarder directement la
mémoire de l'optimiseur en plus des poids dans les futurs instantanés
(\texttt{-{}-checkpoint-optimizer-state}), pour éviter ce réchauffage dans les prochaines campagnes.
Il n'a pas été activé rétroactivement, car cela obligerait à recommencer les entraînements déjà
réalisés — ce qui, au vu du point (1) ci-dessus, n'apporterait pas d'information nouvelle par
rapport à un simple réchauffage.

% =====================================================================
\section{Amortissement et échelle : comment traiter le paramètre $\lambda$}
% =====================================================================

Toute métrique qui passe par un inverse (impossible à calculer directement sur une matrice qui a des
valeurs propres nulles ou proches de zéro) a besoin d'un petit terme d'amortissement $\lambda$
ajouté sur la diagonale avant l'inversion. Plutôt que de fixer une seule valeur de $\lambda$
(ce qui reviendrait à figer arbitrairement un compromis), la campagne balaie systématiquement
$\lambda=\alpha\,\bar\lambda$, avec $\bar\lambda=\mathrm{tr}(R)/P$ (une échelle naturelle liée à la
taille moyenne des valeurs propres de la référence) et $\alpha\in\{10^{-4},10^{-3},10^{-2},10^{-1},1\}$,
et présente les résultats sous forme de courbes plutôt que de points isolés. La valeur
$\lambda=10^{-3}$ réellement utilisée par AdaFisher en entraînement n'est ajoutée qu'en protocole
P2, car elle n'a de sens que pour des facteurs déjà renormalisés entre 0 et 1 par le min-max — un
traitement propre au mode \texttt{diag} qui n'a pas d'équivalent dans les quatre autres modes.

Pour les approximations qui ne prétendent, par construction, respecter aucune échelle absolue
(la diagonale brute d'AdaFisher, ou le second moment libre d'Adam), on rapporte en plus une version
« recalée au mieux », en multipliant l'approximation par le facteur scalaire
$c^\star=\langle R,K\rangle_F/\|K\|_F^2$ qui minimise l'écart de Frobenius — de façon à ne pas
pénaliser injustement une méthode pour un simple facteur d'échelle qu'elle n'a jamais prétendu
respecter.

% =====================================================================
\section{Le plancher de bruit : à partir de quand une différence est-elle réelle ?}
% =====================================================================

\subsection{Pourquoi c'est nécessaire}

$F$ et $\hat E$ sont elles-mêmes calculées sur un nombre fini $N$ d'exemples sondes : ce sont donc
des estimations, avec leur propre variabilité statistique. Si l'on compare deux approximations et
que l'écart mesuré entre elles est plus petit que la variabilité naturelle de $F$ elle-même
(c'est-à-dire l'écart qu'on obtiendrait simplement en changeant l'échantillon d'exemples utilisés
pour la calculer), cette différence ne veut rien dire : elle pourrait tout aussi bien disparaître,
voire s'inverser, avec un autre tirage d'exemples.

\subsection{Comment on le mesure}

On coupe l'ensemble des exemples sondes en deux moitiés aléatoires, on calcule $F$ séparément sur
chaque moitié, et on mesure l'écart entre les deux résultats — en répétant l'opération sur 20
partitions aléatoires différentes pour obtenir un intervalle statistique à 95\,\%. \textbf{Toute
différence entre deux approximations qui n'excède pas ce plancher n'est simplement pas
interprétable.} C'est un principe de lecture des résultats, pas une hypothèse en soi (dans le
brouillon initial, cette règle porte le nom d'hypothèse HF8 ; dans le plan actuellement suivi, elle
est reclassée comme une \emph{règle de lecture} systématique plutôt que comme une hypothèse à
trancher une fois pour toutes).

\subsection{Sondes d'entraînement contre sondes de validation}

Un deuxième axe de comparaison consiste à utiliser deux ensembles d'exemples sondes différents : des
exemples que le modèle a effectivement vus pendant son entraînement, et des exemples qu'il n'a
jamais vus (un ensemble de validation, tiré de la même partition aléatoire que celle réellement
utilisée par l'entraînement, pour garantir que « entraînement » désigne bien les points sur lesquels
le modèle a réellement appris). L'écart entre $F$ et $\hat E$ pourrait très bien se comporter
différemment sur ces deux ensembles — c'est exactement ce que teste l'hypothèse HF1, détaillée
section~\ref{sec:hypotheses}.

% =====================================================================
\section{Les points de contrôle, les « bras » de comparaison, et les graines aléatoires}
\label{sec:checkpoints}
% =====================================================================

\subsection{Une trajectoire d'entraînement, cinq instantanés}

Pour chaque modèle, les bancs d'essai de ce dépôt sauvegardent déjà, à des fractions fixes de
l'entraînement complet (0\,\%, 1\,\%, 10\,\%, 50\,\% et 100\,\%), un instantané des poids du réseau.
Ces fractions sont calculées relativement à une trajectoire de référence \textbf{commune}, partagée
par toutes les stratégies d'entraînement comparées (ce que le plan appelle des « bras » : sept au
total — les cinq modes de courbure, plus Adam, plus AdamW), de sorte qu'un instantané « à 50\,\% »
représente bien la même \emph{quantité d'entraînement} quel que soit le bras, même si les
différentes stratégies n'avancent pas forcément à la même vitesse en temps réel. C'est cette
trajectoire commune, avec ses cinq instantanés, qui constitue l'axe principal le long duquel la
campagne mesure comment la fidélité des approximations \emph{évolue} au fil de l'entraînement
(plutôt qu'une mesure ponctuelle à un seul instant, comme le fait l'article original d'AdaFisher).

\subsection{Un piège déjà rencontré et corrigé}

Une subtilité pratique, découverte lors d'une première campagne d'entraînement sur le cluster de
calcul : pour un bras qui ne dispose que d'un budget de temps fixe plutôt que d'un nombre d'époques
fixe, la fraction « réellement atteinte » à la fin de son propre entraînement peut être inférieure à
100\,\%. Un instantané marqué comme « non programmé » signale que ce n'est pas la fraction nominale
de 100\,\% qui a été atteinte, mais simplement la fin propre à ce bras — une distinction essentielle
pour ne pas comparer par erreur des quantités d'entraînement différentes sous une même étiquette.

\subsection{L'état des graines aléatoires}

Chaque expérience est en principe répétée sur plusieurs graines aléatoires différentes (des
initialisations et des ordres de mini-batchs différents), pour s'assurer qu'une conclusion ne tient
pas à un seul tirage de chance. Au moment de la rédaction, la graine 0 est disponible pour les sept
bras sur six des huit modèles ; des graines supplémentaires sont en cours d'exécution sur le
cluster de calcul pour un sous-ensemble de bras (uniquement \texttt{diag} et \texttt{adamw}, les
deux stratégies jugées les plus informatives pour cet axe). La règle de lecture explicite est que
\textbf{tout résultat agrégé sur plusieurs graines doit annoncer le nombre de graines réellement
disponibles au moment de la mesure}, plutôt que de supposer par avance un nombre cible qui n'est pas
encore garanti.

% =====================================================================
\section{Les métriques, une par une}
\label{sec:metriques}
% =====================================================================

Toutes les métriques ci-dessous comparent une référence exacte $R$ (soit $F$, soit $\hat E$) à une
approximation $K$, avec $X_\lambda=X+\lambda I$ pour désigner une matrice amortie. Elles sont
choisies pour capturer des aspects \textbf{complémentaires} de la ressemblance entre deux matrices :
aucune métrique seule ne suffirait, comme l'illustre l'exemple suivant. Deux matrices peuvent avoir
un écart de Frobenius très faible parce qu'elles se ressemblent presque partout, tout en étant
radicalement différentes dans une seule direction très importante — c'est exactement ce que les
métriques ci-dessous (M1 à M8) sont chacune conçues pour détecter ou, au contraire, pour ignorer.

\subsection{M1 — l'écart entrée par entrée (Frobenius)}

\[
e_F = \frac{\|R-K\|_F}{\|R\|_F}, \qquad
\cos_F = \frac{\langle R,K\rangle_F}{\|R\|_F\,\|K\|_F}, \qquad
e_F^\star=\sqrt{1-\cos_F^2}.
\]
C'est la mesure la plus simple : on traite les deux matrices comme deux tableaux de nombres et on
mesure leur écart quadratique moyen, normalisé par la taille de la référence — l'équivalent de
comparer deux images pixel par pixel. Le produit scalaire de Frobenius $\langle R,K\rangle_F$ donne
un « cosinus » entre les deux matrices, qui ignore leur échelle relative ($e_F^\star$, la version
« à échelle optimale », retire ce facteur d'échelle avant de mesurer l'écart résiduel). Cette
métrique a un défaut connu : elle est \textbf{dominée} par les plus grandes valeurs propres, et peut
donc masquer une erreur sévère concentrée dans une direction de faible amplitude mais importante
pour l'optimisation.

\subsection{M2 — la pire direction (norme spectrale)}

\[
e_2 = \frac{\|R - c^\star K\|_2}{\|R\|_2}.
\]
Alors que M1 moyenne l'erreur sur toutes les directions, M2 ne regarde que la \textbf{direction où
l'erreur est la plus grande} (la norme spectrale, ou norme d'opérateur, d'une matrice symétrique est
sa plus grande valeur propre en valeur absolue). C'est une mesure de robustesse dans le pire des
cas, complémentaire de M1.

\subsection{M3 — la divergence de Kullback--Leibler entre deux gaussiennes (perte de Stein)}

\[
D_\lambda(K\|R) = \tfrac12\Big[\mathrm{tr}(K_\lambda R_\lambda^{-1}) - P - \log\det(K_\lambda R_\lambda^{-1})\Big].
\]
Cette métrique demande un peu plus de mise en contexte. Toute matrice de courbure amortie $R_\lambda$
peut être vue comme la matrice de covariance \emph{inverse} d'une distribution gaussienne
hypothétique sur les déplacements de paramètres — c'est exactement l'objet mathématique que le
gradient naturel utilise implicitement pour définir « la bonne direction, à la bonne échelle ».
$D_\lambda(K\|R)$ est alors, tout simplement, la divergence de Kullback--Leibler classique entre les
deux distributions gaussiennes associées à $K_\lambda$ et $R_\lambda$ — une façon standard, en
théorie de l'information, de mesurer à quel point deux distributions de probabilité sont
différentes.

L'intérêt décisif de cette métrique par rapport à M1/M2 est qu'elle est \textbf{invariante par
reparamétrage linéaire} : si l'on change arbitrairement d'échelle ou de base de coordonnées pour les
paramètres (un changement qui ne devrait rien changer au \emph{comportement} d'un préconditionneur,
seulement à sa représentation numérique), $D_\lambda$ reste identique, alors que M1 et M2 changeraient.
C'est donc le critère le plus fidèle à ce qui compte vraiment pour un préconditionneur : non pas « à
quel point les nombres se ressemblent », mais « à quel point les deux définissent la même notion de
distance dans l'espace des paramètres ». Un test de cohérence trivial mais rassurant : quand $K=R$,
cette divergence vaut exactement zéro.

\subsection{M4 — le conditionnement (vitesse de convergence prévisible)}

\[
\kappa_\lambda = \frac{\lambda_{\max}}{\lambda_{\min}} \text{ de } K_\lambda^{-1/2}R_\lambda K_\lambda^{-1/2}.
\]
Cette métrique répond à une question très pratique : si l'on utilisait $K$ comme préconditionneur
dans une méthode itérative de type gradient conjugué pour résoudre un problème dont la vraie
courbure est $R$, combien d'itérations faudrait-il ? La théorie classique de l'optimisation dit que
ce nombre d'itérations croît comme $\sqrt{\kappa}$ (la racine carrée du rapport entre la plus grande
et la plus petite valeur propre de l'opérateur « relatif » ci-dessus) : un $\kappa$ proche de 1
signifie que $K$ est un préconditionneur presque parfait pour $R$.

\subsection{M5 — l'efficacité du pas naturel ($\rho$), la métrique la plus directement pratique}

\[
\rho(K) = \frac{(g^\top d)^2}{(d^\top R_\lambda d)(g^\top R_\lambda^{-1}g)}, \qquad d = K_\lambda^{-1}g.
\]
C'est probablement la métrique la plus intuitive et la plus directement liée à l'utilisation réelle
d'une approximation de courbure en entraînement. Étant donné un gradient $g$, on calcule le pas
$d=K_\lambda^{-1}g$ que l'approximation \emph{proposerait} de suivre. $\rho(K)$ compare la baisse de
perte que ce pas $d$ permettrait effectivement d'obtenir (dans un modèle quadratique local dont la
vraie courbure est $R$) à la baisse de perte \textbf{optimale} qu'on obtiendrait avec la meilleure
direction possible, $R_\lambda^{-1}g$. Par une inégalité de Cauchy--Schwarz, $\rho$ est toujours
compris entre 0 et 1 : $\rho=1$ signifie que l'approximation, même compressée, donne exactement le
même pas que la référence exacte ; $\rho$ proche de 0 signifie que le pas proposé est presque
inutile, voire contre-productif.

C'est cette métrique que la section~\ref{sec:stall} identifie comme la mesure la plus directement
décisive pour comprendre l'échec observé du banc principal (l'auto-encodeur MNIST) : elle permet de
demander directement, à l'endroit précis où l'entraînement se bloque, « le pas que l'optimiseur
propose réellement est-il, ou non, encore une bonne approximation du pas idéal ? ».

\subsection{M6 — le recouvrement des directions dominantes}

\[
\frac{\|V_R^\top V_K\|_F^2}{k}, \qquad k\in\{C,\,10C\}.
\]
Les spectres de courbure des réseaux de neurones ont une propriété bien documentée dans la
littérature : un petit nombre de directions (de l'ordre du nombre de classes $C$) concentrent
l'essentiel de l'énergie, le reste étant beaucoup plus plat. M6 mesure si les quelques directions les
plus importantes de l'approximation ($V_K$, ses vecteurs propres dominants) coïncident bien avec les
quelques directions les plus importantes de la vraie référence ($V_R$) — indépendamment de ce qui se
passe dans le reste, beaucoup plus plat, du spectre.

\subsection{M7 — le biais d'indépendance et la part du partage de poids}

Cette famille de mesures quantifie directement, plutôt que de le supposer, à quel point
l'hypothèse d'indépendance sur laquelle repose K-FAC est vraie sur les réseaux réels : le rapport
entre la deuxième et la première valeur singulière (proche de 0 si le bloc exact ressemble
vraiment à un simple produit de Kronecker, plus grand sinon) et la part d'énergie totale qui n'est
pas capturée par cette meilleure approximation en produit de Kronecker possible, obtenue par
décomposition en valeurs singulières (SVD) d'un réarrangement particulier de la matrice
$B_\ell^{\rm expand}$. Le même outil sert à isoler, séparément, la part perdue à cause du partage de
poids ($\|B_\ell - B_\ell^{\rm expand}\|_F/\|B_\ell\|_F$) et la part perdue en ne gardant que la
diagonale du produit de Kronecker.

\subsection{M8 — le couplage entre couches différentes}

\[
c_{\ell\ell'} = \frac{\langle K_\ell, K_{\ell'}\rangle_F}{\|K_\ell\|_F\,\|K_{\ell'}\|_F}.
\]
Cette métrique quantifie la validité de l'approximation la plus radicale de toutes citées plus
haut : le bloc-diagonal exact, qui suppose qu'aucune information de courbure ne circule entre deux
couches différentes. C'est mathématiquement l'équivalent d'une mesure de similarité (une forme non
centrée de ce qu'on appelle en apprentissage automatique une « CKA », \eng{centered kernel
alignment}) entre les « noyaux tangents » de deux couches — un $c_{\ell\ell'}$ proche de zéro
confirmerait que les couches sont bien quasiment indépendantes en termes de courbure ; une valeur
substantielle indiquerait, au contraire, une information réelle perdue par l'approximation
bloc-diagonale.

\subsection{Comment calculer M3 et M5 sans jamais former de matrice $P\times P$ (régime B)}

Un point technique important, qui explique pourquoi le régime B (streamé, couche par couche) peut
malgré tout calculer les métriques M3 (KL) et M5 ($\rho$), pourtant définies à partir d'inverses de
matrices $P\times P$ : en réécrivant $R=U^\top U$ avec $U\in\mathbb R^{m\times P}$ ($m\ll P$
possible), l'identité de Woodbury permet de remplacer chaque inverse $P\times P$ par un inverse
$m\times m$, bien moins coûteux :
\[
R_\lambda^{-1} = \lambda^{-1}\big(I - U^\top M^{-1} U\big), \qquad M = \lambda I_m + UU^\top,
\]
et de même pour le log-déterminant, via l'identité de Sylvester,
$\log\det R_\lambda = P\log\lambda + \log\det M - m\log\lambda$. C'est le même principe
mathématique qui rend les filtres de Kalman efficaces, appliqué ici à une matrice de courbure qui
n'est, par construction, jamais que le produit d'un nombre fini d'exemples sondes.

% =====================================================================
\section{Les hypothèses et questions, remises en contexte}
\label{sec:hypotheses}
% =====================================================================

Le brouillon original du plan (v0) formulait huit hypothèses \textbf{falsifiables} — chacune
accompagnée explicitement de la condition qui la réfuterait, une exigence de rigueur scientifique
délibérée. Le plan actuellement suivi par le code (v1) resserre l'ordre de priorité de ces huit
hypothèses (certaines sont traitées en premier, d'autres reportées, une devient une simple règle de
lecture plutôt qu'une question ouverte) ; les deux niveaux sont présentés ici ensemble, avec leur
statut actuel.

\subsection*{HF1 — l'écart source-vs-vraie-Fisher se dégrade-t-il différemment sur l'entraînement et sur la validation ?}
\textit{Énoncé :} sur des exemples tirés de l'ensemble d'entraînement, l'écart entre la Fisher
empirique et la vraie Fisher (mesuré par l'efficacité $\rho$) se dégrade au fil de l'entraînement ;
sur des exemples de validation jamais vus, beaucoup moins. \textit{Intuition :} sur des points déjà
appris, le gradient d'un exemple individuel tend vers zéro d'une façon particulière qui n'a pas de
raison de se produire sur des points non vus — c'est une forme du même phénomène de surapprentissage
qui affecte la perte elle-même, mais appliqué ici à la courbure plutôt qu'à la valeur de la perte.
\textit{Réfutée si} les deux courbes (entraînement, validation) coïncident à l'intérieur du plancher
de bruit. \textit{Statut actuel :} priorisée dans v1 ; une première mesure sur le modèle A1 avait laissé un
résultat ambigu entre les trois combinaisons possibles (entraînement/validation, entraînement/test,
validation/test), faute d'un dénominateur commun enregistré pour les trois. Ce défaut
d'instrumentation est corrigé — chaque écart est maintenant rapporté avec les deux normes qui le
divisent — et la cause de l'ambiguïté est identifiée : la courbure mesurée sur des exemples jamais
vus est réellement \emph{plus grande} (de 30 à 60\,\% sur A1), parce qu'à mi-parcours le modèle est
plus confiant sur ce qu'il a appris. L'hypothèse elle-même reste ouverte : aucun des calculs menés
jusqu'ici n'a utilisé de sondes de validation, tous ayant porté sur des sondes d'entraînement.

\subsection*{HF2 — l'erreur de structure de K-FAC dépend-elle du type de couche, et \eng{expand} contre \eng{reduce} ?}
\textit{Énoncé :} l'erreur de structure de K-FAC est plus grande sur les couches à poids partagés
(attention, convolution) que sur les couches simples, et le choix \eng{reduce} (qui résume d'abord
chaque exemple par une moyenne sur ses positions avant de corréler) est préférable en classification
d'image, l'inverse en modélisation du langage. \textit{Réfutée si} l'écart entre les deux choix est
sous le plancher de bruit. \textit{Statut :} \textbf{tranchée — confirmée sur les trois modèles à
poids partagés} (section~\ref{sec:lot3}). Les deux volets le sont : l'erreur de structure est bien
plus grande sur les couches à poids partagés que sur la tête de classification du même modèle, et
\eng{reduce} l'emporte sur \eng{expand} dans 14/15, 15/15 et 40/45 des cas, très au-dessus du
plancher de bruit. Deux réserves accompagnent ce verdict, toutes deux mesurées : l'erreur ainsi
mesurée provient \emph{d'abord} du partage de poids lui-même (97\,\% de l'écart) et non de
l'hypothèse d'indépendance ; et le classement s'inverse si l'on juge non plus la fidélité mais la
qualité du pas produit ($\rho$), ce qui arrive dans environ 70\,\% des cas. Le volet « l'inverse
en modélisation du langage » reste, lui, non testé : aucun modèle de langage n'est encore construit
dans ce dépôt.

\subsection*{HF3 — l'hypothèse d'indépendance d'AdaFisher a-t-elle un coût mesurable sur la seule diagonale ?}
\textit{Énoncé :} l'écart entre les deux chemins vers la diagonale d'AdaFisher (décrit
section~4.6, une covariance $\mathrm{Cov}(a_j^2,g_i^2)$) dépasse 10\,\% sur les couches de
normalisation et sur l'attention. \textit{Réfutée si} les deux chemins coïncident à 5\,\% près.
\textit{Statut :} \textbf{tranchée — confirmée sur les quatre modèles déjà traités}
(sections~\ref{sec:lot2} et~\ref{sec:lot3}). Le raccourci coûte $0{,}86$--$0{,}90$ sur les
convolutions et $0{,}83$--$0{,}97$ sur les couches d'attention et de MLP du transformeur en fin
d'entraînement, contre $0{,}25$--$0{,}41$ sur les têtes de classification : un ordre de grandeur
au-delà du seuil de 10\,\% que l'hypothèse s'était fixé. La seule exception connue est la première
couche du modèle A1, et elle s'explique entièrement par les pixels toujours nuls de MNIST.

\subsection*{HF4 — le couplage $\gamma$--$\beta$ ignoré par AdaFisher dans les normalisations est-il négligeable ?}
\textit{Énoncé :} dans une couche de normalisation, le terme croisé entre les deux paramètres
$\gamma$ et $\beta$ — que la proposition 3.1 de l'article AdaFisher exclut explicitement — porte une
fraction non négligeable de l'énergie du bloc exact complet. \textit{Statut :} déjà en grande partie
résolue \textbf{analytiquement}, au moment même de construire les cinq modes de l'optimiseur (avant
cette campagne) : le bloc exact d'une normalisation a une structure de Hadamard, pas de Kronecker, et
le petit terme croisé y a été délibérément conservé plutôt que forcé à zéro. Cette mesure est
maintenant faite sur les quatre modèles : le terme croisé porte environ deux tiers de l'énergie du
bloc exact sur le modèle A1, $0{,}65$--$0{,}70$ sur le convolutif et $0{,}39$--$0{,}62$ sur le
transformeur — non négligeable partout, donc.

\textit{Un résultat imprévu du même lot} (section~\ref{sec:lot3}) : la forme de Hadamard décrit
correctement la \emph{direction} de la courbure d'une normalisation, mais en sous-estime la
\emph{taille} d'un facteur 20 à 38 dès qu'un poids est appliqué à plusieurs positions (contre un
facteur $0{,}95$, c'est-à-dire aucun écart, sur le modèle sans partage). Un estimateur corrigé pour
les normalisations doit donc corriger l'échelle, et pas seulement réintroduire le terme croisé. Et
la formule effectivement livrée aujourd'hui dans le code se classe dernière des cinq lectures
comparées, y compris une fois l'échelle neutralisée.

\subsection*{HF5 — l'information qui circule entre deux couches différentes compte-t-elle vraiment ?}
\textit{Énoncé :} ignorer entièrement les corrélations entre couches (l'approximation
bloc-diagonale) ne coûte presque rien, aussi bien sur un CNN que sur un transformeur — sauf peut-être
sur le couplage spécifique entre les projections \eng{query} et \eng{key} de l'attention. \textit{Réfutée
si} l'écart est grand sur le transformeur. \textit{Statut :} première réponse partielle, et elle va
dans le sens de l'hypothèse — mais dans une direction inattendue. Le couplage mesuré entre deux
couches (section~\ref{sec:lot3}) est \emph{plus faible} là où les poids sont partagés : au minimum
$0{,}28$ sur le transformeur et $0{,}44$ sur le convolutif, contre $0{,}63$ sur le perceptron A1.
Ignorer les corrélations inter-couches est donc mieux justifié sur les architectures réalistes que
sur le modèle jouet où la question avait d'abord été posée. Le verdict complet demande encore la
mesure directe de $\rho$ pour l'approximation bloc-diagonale, et la séparation \eng{query}/\eng{key}
que les modèles actuels ne permettent pas (leur projection est fusionnée).

\subsection*{HF6 — un tirage Monte-Carlo unique dégrade-t-il l'estimation de courbure ?}
\textit{Énoncé :} avec un seul tirage aléatoire ($K=1$) et un petit \eng{batch}, l'estimation de type
Monte-Carlo est moins bonne (en $\rho$) que l'estimation purement empirique — un résultat déjà
suggéré par des travaux antérieurs sur la Fisher empirique inversée. \textit{Statut :} repoussée dans
l'ordre de priorité de v1 (elle n'est testée qu'à travers les sources Monte-Carlo optionnelles,
$K\in\{1,4,16\}$, qui ne sont pas la priorité des premiers lots).

\subsection*{HF7 — le classement « en pratique » (P2) est-il le même que le classement « dans l'idéal » (P1) ?}
\textit{Énoncé :} le classement des cinq approximations, mesuré dans les conditions réelles
d'entraînement (P2), n'est \textbf{pas} le même que celui mesuré dans les conditions idéalisées (P1)
— autrement dit, ce sont le lissage temporel et la renormalisation, pas la structure mathématique
choisie, qui dominent l'erreur réellement subie par AdaFisher en pratique. \textit{Statut :}
priorisée dans v1 ; les mesures préliminaires de fidélité du réchauffage (section~7.4) vont déjà dans
ce sens en creux — un réchauffage insuffisant peut, à lui seul, fausser le préconditionneur appliqué
de plusieurs dizaines de pourcents, indépendamment de toute question de structure.

\subsection*{HF8 — une partie des différences mesurées est-elle simplement du bruit ?}
\textit{Énoncé :} à petite taille d'échantillon, une partie des différences apparentes entre
approximations est en réalité indiscernable de la variabilité statistique propre à $F$ elle-même.
\textit{Statut dans v1 :} cette hypothèse n'est plus traitée comme une question à trancher une fois,
mais promue en \textbf{règle de lecture systématique} appliquée à \emph{tous} les résultats de la
campagne (section~\ref{sec:hypotheses} plus haut, le « plancher de bruit ») — un changement de
statut logique plutôt qu'un abandon. \textit{Constat pratique} : ce plancher dépend énormément du
modèle et du nombre d'exemples sondes. Il valait $0{,}343$ sur le modèle A1 (55\,000 exemples), au
point de rendre inexploitable un premier essai à 512 exemples ; il tombe à $0{,}005$--$0{,}039$ sur
les trois modèles du lot 3 (45\,000 exemples), où les écarts mesurés lui sont dix à cent fois
supérieurs. La règle reste indispensable — c'est elle qui avait invalidé l'essai à 512 exemples —
mais elle ne mord pas sur les résultats du lot 3.

\subsection*{Les questions Q1 à Q5}

Les questions Q1 à Q4 (déjà détaillées section~\ref{sec:design}) structurent le design factoriel
lui-même — elles ne sont pas des hypothèses au sens strict (elles ne prédisent rien à l'avance),
mais des axes de décomposition de l'erreur totale. La question \textbf{Q5} du brouillon original
(« que devient le classement en conditions opérationnelles réelles ? ») est exactement ce que le
protocole P2 (section~\ref{sec:protocoles}) répond directement dans le plan actuel ; elle n'est donc
plus listée séparément, ayant simplement été absorbée dans la définition même de P2.

% =====================================================================
\section{Ce qui a déjà été mesuré}
\label{sec:mesures}
% =====================================================================

Cette section rassemble, en un seul endroit, tout ce que la campagne a déjà produit comme résultats
concrets — par opposition au reste de ce document, qui décrit surtout ce qui est \emph{prévu}.

\subsection{Le premier calcul dense complet (modèle A1, sur la totalité des 55\,000 images d'entraînement de MNIST)}

Un premier calcul complet et exact de $F$ et $\hat E$ (sans aucune structure imposée) a été mené sur
le modèle A1, en utilisant \textbf{tout} l'ensemble d'entraînement de MNIST comme exemples sondes
(55\,000 images), sur une durée d'un peu plus d'une heure de calcul GPU. Plusieurs constats mesurés
en sont sortis :

\begin{mesure}[Résultats mesurés — le calcul dense A1]
\begin{itemize}[leftmargin=1.3em]
\item Le plancher de bruit propre à $F$, estimé en coupant les 55\,000 exemples en deux moitiés,
vaut $0{,}2986$ (soit un écart-type de bruit $\sigma_N=0{,}1493$). L'écart mesuré entre la Fisher
vraie et la Fisher empirique (Q1) vaut $0{,}3751$, soit \textbf{2,51 fois} ce plancher de bruit :
l'écart source-Fisher est donc réellement interprétable sur ce modèle, et non un artefact
d'échantillonnage.
\item Ce plancher de bruit décroît plus lentement que ce que prédit la théorie de l'estimation
statistique classique : en passant de 4\,000 à 55\,000 exemples (un facteur 13,75), le bruit n'a
diminué que d'un facteur 2,92, alors qu'un estimateur « idéal » aurait dû s'améliorer d'un facteur
3,71 ($\propto N^{-1/2}$). Autrement dit, MNIST n'offre plus beaucoup de marge d'amélioration
supplémentaire par ce seul levier.
\item Le rang réel de $F$ (21\,829 sur 26\,634 paramètres possibles) et celui de $\hat E$ (18\,342)
montrent que, une fois le nombre d'exemples suffisamment grand, c'est la Fisher \emph{empirique} qui
est structurellement la plus déficiente en rang — un résultat plus net que ce que la seule théorie
générale (Kunstner et al., 2019) permettait d'anticiper.
\item Chaque déficience de rang a une explication structurelle précise : le bloc de la tête de
classification a un noyau de dimension exactement $d_{\rm in}+1$ (ajouter une même constante à
chaque logit ne change rien à la sortie du softmax — un fait mathématique, pas un artefact
numérique) ; et le premier bloc du réseau a un noyau de dimension 4\,579, presque exactement
$32\times143$, qui correspond aux pixels de bordure de MNIST, identiquement nuls sur toutes les
images.
\item Ce même premier bloc du réseau concentre à lui seul \textbf{63\,\%} de l'énergie totale
(la trace) de $F$ : sur ce modèle particulier, la courbure du problème d'apprentissage est
massivement dominée par sa toute première couche, un fait qui doit conditionner l'interprétation de
toute conclusion « par type de couche » tirée de ce modèle.
\end{itemize}
\end{mesure}

Ce même calcul a également mis au jour, en amont, un piège de conception facile à manquer :
l'accroche technique utilisée par défaut pour capturer les gradients (\texttt{register\_full\_
backward\_hook}) ne se déclenche \textbf{pas} pour la toute première couche du réseau, dès lors que
l'entrée du réseau elle-même ne demande pas de gradient — ce qui aurait silencieusement mis à zéro
94\,\% des paramètres du modèle A1 dans le calcul de référence, sans qu'aucune erreur ne soit
levée. Le code de la campagne utilise donc une méthode différente et plus sûre pour capturer ces
gradients, avec un test de non-régression dédié qui vérifie explicitement ce cas.

\subsection{Le blocage de l'auto-encodeur MNIST}
\label{sec:stall}

Le banc d'essai \textbf{principal} du projet (l'auto-encodeur MNIST, huit couches linéaires avec des
sigmoïdes) présente une anomalie déjà identifiée avant même le début de cette campagne, et que la
campagne est justement conçue pour aider à élucider :

\begin{mesure}[Constat mesuré — le blocage de l'auto-encodeur]
Les cinq modes de courbure (\texttt{diag}, \texttt{kfac}, \texttt{ekfac}, \texttt{tkfac},
\texttt{tekfac}) se bloquent tous à une erreur quadratique moyenne d'environ $0{,}7085$ dès les
deux premières époques d'entraînement, et n'en bougent quasiment plus ensuite : entre 10\,\% et
100\,\% de la durée totale de l'entraînement, les paramètres ne se déplacent plus que de
$0{,}2$ à $0{,}5\,\%$. Dans les mêmes conditions, Adam continue de progresser normalement jusqu'à
une erreur de $0{,}4836$, en déplaçant ses paramètres de $147\,\%$ sur la même période. Un
\eng{sweep} du terme d'amortissement $\lambda$ sur quatre ordres de grandeur a déjà exclu ce
paramètre comme cause principale : il ne modifie le résultat final que de moins de 1\,\% pour les
quatre modes de type Kronecker.
\end{mesure}

Les deux causes restantes, non encore testées avant cette campagne, sont le taux d'apprentissage
partagé par tous les modes, et \textbf{la géométrie même du produit de Kronecker $A\otimes B$ sur ce
réseau particulier} — c'est-à-dire la possibilité que l'approximation en produit de Kronecker soit,
sur ce réseau précis, structurellement mal adaptée (plutôt que simplement mal réglée). C'est
précisément cette deuxième cause que la métrique M5 ($\rho$, l'efficacité du pas naturel) et le
rapport $\|\tilde F^{-1}\hat m\|/\|\hat m\|$ (à quel point le pas préconditionné diffère en norme du
gradient brut) sont conçus pour trancher directement, en comparant, au même point de blocage, la
trajectoire produite par un mode de type Kronecker à celle produite par Adam. Les instantanés
nécessaires à cette mesure existent déjà (ils ont été sauvegardés lors des campagnes
d'entraînement précédentes) ; c'est la mesure la plus rapide et la plus décisive que le premier
travail en régime B de cette campagne doit produire, précisément parce qu'elle répond à une question
que le développement de l'optimiseur, à lui seul, ne peut actuellement pas trancher.

\subsection{Le zoo complet des approximations sur le modèle A1 (lot 2)}
\label{sec:lot2}

Le deuxième lot a appliqué au modèle A1 l'ensemble des approximations décrites en
section~\ref{sec:zoo} et les métriques M1, M3, M5, M7, M8, à ses cinq instantanés d'entraînement et
pour les deux sources de signal. Trois choses en sont ressorties.

D'abord, \textbf{le raccourci diagonal d'AdaFisher n'est pas gratuit} (HF3) : sur la toute première
couche du réseau il ne coûte presque rien — mais pour une raison très particulière, à savoir que
130 des 784 pixels de MNIST sont identiquement nuls, et qu'une quantité toujours nulle ne peut
corréler avec rien ; sur toutes les couches situées en aval, l'écart est réel et \emph{croît} au fil
de l'entraînement.

Ensuite, un résultat inattendu, qui n'était pas la question posée : \textbf{toutes les
approximations se dégradent nettement au fil de l'entraînement} quand on les juge sur la qualité du
pas qu'elles produiraient (métrique $\rho$). En début d'entraînement, les approximations structurées
récupèrent 61 à 97\,\% du pas idéal ; à la fin, seulement 13 à 46\,\%, et les approximations
diagonales tombent à une fraction de pourcent sur les premières couches. Une approximation presque
exacte au départ est un mauvais substitut de la vérité à l'arrivée.

Enfin, une \textbf{correction apportée après coup} à la lecture de ce lot, et qu'il faut connaître
avant de lire la suite. Sur les couches de normalisation, quatre lectures simplifiées avaient été
comparées, et le classement obtenu avait été présenté comme un résultat. Le lot 3 a établi que
trois des quatre comparaisons de ce classement sont en réalité des \textbf{théorèmes} — vraies
d'avance, quelles que soient les données : la diagonale exacte est, par construction, la meilleure
approximation diagonale possible au sens de Frobenius, et le bloc « séparé » est la meilleure
matrice sans termes croisés. Seule la comparaison entre la forme de Hadamard et la diagonale exacte
était une véritable mesure. De plus, la lecture présentée comme « la formule réellement utilisée par
l'optimiseur » n'était pas cette formule, mais la diagonale de la forme de Hadamard. Les deux points
sont corrigés dans le lot 3, qui calcule la formule réellement livrée en appelant directement le
code de l'optimiseur, et qui vérifie désormais les trois inégalités-théorèmes à l'exécution : leur
violation serait un bogue, jamais un résultat.

\subsection{Les modèles à poids partagés : convolutions et attention (lot 3)}
\label{sec:lot3}

Le troisième lot étend tout ce qui précède aux deux modèles où un même poids est appliqué à
\emph{plusieurs endroits} : le petit réseau convolutif A2 (chaque filtre balaie l'image, soit 1\,024,
256 puis 64 positions selon la couche), dans ses deux variantes de normalisation, et le petit
transformeur A3 (la même couche traite les 64 jetons de l'image). C'est la situation pour laquelle
K-FAC a été généralisé, et celle que le modèle A1 — où chaque poids ne sert qu'une fois — ne pouvait
pas tester du tout.

\subsubsection*{Avant de mesurer : trois erreurs dans du code jamais exécuté}

Le code du lot 2 traitait déjà, sur le papier, les couches à poids partagés ; simplement, aucun
modèle ne l'avait jamais fait tourner. Écrire d'abord les tests, puis les lancer contre ce code, a
révélé trois erreurs, chacune mesurée plutôt que soupçonnée :

\begin{mesure}[Erreurs trouvées et corrigées avant toute mesure]
\begin{itemize}[leftmargin=1.3em]
\item La variante \eng{reduce} de K-FAC sommait les entrées sur les positions là où la définition
publiée en prend la \emph{moyenne} : le résultat était exactement $T^2$ fois trop grand ($T$ = nombre
de positions). Mesuré sur deux cas-jouets : erreur relative de $24{,}000$ pour $T=5$ et de
$80{,}000$ pour $T=9$, soit précisément $T^2-1$. Le biais allait dans le sens qui \emph{avantage}
l'autre variante — c'est-à-dire exactement la comparaison que ce lot devait trancher.
\item La version « à valeurs propres corrigées » de cette même variante projetait une quantité qui
n'est le gradient de rien, ce qui lui faisait perdre la propriété de conservation d'énergie qui la
définit.
\item TKFAC, sur ces couches, utilisait une troisième formule ne conservant l'énergie d'aucun objet
défini, et différente de celle que l'optimiseur du dépôt applique réellement.
\end{itemize}
Les tests correspondants échouent sur l'ancien code et passent sur le nouveau. Aucun résultat
antérieur n'est affecté : ces trois branches ne s'activent que lorsqu'un poids est partagé, ce qui
n'arrive jamais sur le modèle A1.
\end{mesure}

Deux autres défauts ont été corrigés au passage. Sur le transformeur, 2\,048 des 21\,098 paramètres
(les positions apprises) étaient absents de la matrice de référence, laquelle était malgré tout
étiquetée « la vraie courbure » ; ils sont désormais inclus. Et chaque construction de référence se
vérifie maintenant elle-même contre le calcul automatique de gradients, de deux façons
complémentaires — la somme sur tous les exemples, et un exemple isolé — la seconde attrapant ce que
la première ne peut pas voir. Sur les trois modèles réels, l'accord est de l'ordre de
$5\times10^{-16}$, c'est-à-dire la précision de la machine.

\subsubsection*{Ce que les trois calculs ont produit}

Trois calculs GPU, tous terminés sans incident (1\,h\,02 pour le transformeur, 4\,h\,21 et 4\,h\,31
pour les deux variantes du convolutif), chacun sur les \textbf{45\,000} images d'entraînement, aux
cinq instantanés et pour les deux sources. Point important pour la lecture : à cette taille
d'échantillon, le plancher de bruit tombe à $0{,}005$--$0{,}039$ selon le modèle et l'instantané,
contre $0{,}343$ sur le modèle A1. Les écarts rapportés ci-dessous lui sont dix à cent fois
supérieurs — la règle de lecture « en dessous du bruit, on ne conclut pas » ne mord donc plus ici.

Les règles de décision avaient été \textbf{écrites et figées avant de voir la moindre donnée},
précisément pour qu'aucun seuil ne puisse être ajusté après coup.

\begin{mesure}[HF2 tranchée, dans le même sens sur les trois modèles]
\textbf{(ii) Le choix \eng{reduce} est le plus fidèle.} En comptant chaque couple (couche partagée,
instantané) comme un vote, avec son intervalle de confiance :

\begin{center}\small
\begin{tabular}{@{}lcccl@{}}
\toprule
Modèle & cas & \eng{reduce} & \eng{expand} & amplitude des victoires de \eng{reduce} \\
\midrule
A2 (normalisation par groupes) & 15 & \textbf{14} & 1 & $-0{,}71$ à $-0{,}15$ \\
A2 (BatchNorm) & 15 & \textbf{15} & 0 & $-0{,}72$ à $-0{,}11$ \\
A3 (transformeur) & 45 & \textbf{40} & 5 & $-0{,}94$ à $-0{,}08$ \\
\bottomrule
\end{tabular}
\end{center}

\noindent Les rares victoires de \eng{expand} sont petites ($+0{,}02$ à $+0{,}18$) et concentrées
sur les instantanés les plus précoces, à une exception près : la couche qui découpe l'image en
morceaux, à l'initialisation seulement ($+0{,}78$) — la seule couche dont les « positions » sont des
morceaux d'image disjoints plutôt qu'une séquence que le réseau agrège ensuite.

\medskip
\textbf{(i) L'erreur de structure est bien plus grande sur les couches à poids partagés.} L'erreur
médiane de K-FAC y vaut $0{,}95$ à $0{,}98$, contre $0{,}005$--$0{,}011$ à l'initialisation et
$0{,}25$--$0{,}64$ en fin d'entraînement sur la tête de classification, qui n'a pas de partage.
\end{mesure}

\textbf{Mais le résultat le plus important n'était pas la question posée.} La décomposition décrite
en section~\ref{sec:partage} sépare l'erreur de K-FAC en deux étapes : ignorer les corrélations entre positions
différentes, puis supposer l'indépendance entre entrées et gradients. Mesure : la première étape
coûte à elle seule une fraction médiane de \textbf{0,971 — identique sur les trois modèles} (étendue
$0{,}78$ à $0{,}995$), la seconde $0{,}45$ à $0{,}60$ de ce qui reste. Autrement dit, sur une couche
à poids partagés, l'erreur de K-FAC est \textbf{d'abord une erreur de partage de poids}, et non
l'hypothèse d'indépendance dont la littérature discute habituellement.

Un fait complémentaire pointe dans la même direction : à l'initialisation, le bloc exact est
\emph{presque} de la forme que K-FAC suppose (la meilleure approximation de cette forme atteint
$0{,}05$ sur l'attention du transformeur et $0{,}22$--$0{,}37$ sur les convolutions), alors que
K-FAC lui-même reste à $0{,}95$--$0{,}97$. Ce n'est donc pas la forme de Kronecker qui échoue, mais
le choix particulier des deux facteurs que fait K-FAC.

\begin{encadre}[Un avertissement à emporter avec HF2]
Les deux façons de juger une approximation ne donnent pas le même gagnant. Sur l'ensemble des
amortissements et des couches partagées, le choix \eng{reduce} — plus fidèle en norme — produit un
\emph{moins bon} pas d'optimisation que \eng{expand} dans 70/75, 62/75 et 162/225 des cas selon le
modèle. Exemple chiffré (attention du transformeur, fin d'entraînement, amortissement léger) :
erreur $0{,}73$ contre $0{,}97$ en faveur de \eng{reduce}, mais efficacité du pas $\rho = 0{,}14$
contre $0{,}45$ en faveur de \eng{expand}. La règle de décision votait sur la fidélité et
rapportait $\rho$ séparément, précisément parce que les deux pouvaient diverger ; elles divergent.
\textbf{Quiconque choisit entre ces deux variantes pour un optimiseur doit lire $\rho$, pas la
fidélité.}
\end{encadre}

\subsubsection*{Les normalisations : la proposition 3.1 pointe dans la bonne direction, mais 20 à 38 fois trop bas}

La comparaison entre la forme de Hadamard (la proposition 3.1 de l'article AdaFisher) et la simple
diagonale exacte s'inverse entre A1 et les trois modèles de ce lot. L'explication est mesurée, et
c'est une pure question d'échelle : le facteur multiplicatif optimal de la forme de Hadamard vaut
$0{,}95$ sur A1 (aucun partage de poids) mais $27{,}8$, $20{,}4$ et $37{,}9$ sur les trois modèles à
poids partagés. En direction — c'est-à-dire une fois ce facteur d'échelle neutralisé — le classement
est identique sur les quatre modèles. La proposition 3.1 décrit donc correctement la \emph{forme} de
la courbure d'une couche de normalisation, mais en sous-estime la \emph{taille} d'un facteur 20 à 38
dès qu'un poids est appliqué à plusieurs positions, pour une raison structurelle : elle est bâtie
sur des statistiques qui ignorent les corrélations entre positions, et approxime donc le bloc
\emph{déjà amputé} de ces corrélations.

Conséquence directe pour le dépôt : un estimateur corrigé pour les normalisations — piste que le
plan de campagne identifie comme contribution possible — doit corriger l'échelle, et pas seulement
réintroduire le terme croisé $\gamma$--$\beta$. Second point, désormais mesuré sur le bon objet : la
formule effectivement livrée dans le code est la plus faible des cinq lectures testées, y compris en
direction (dernière dans 13 cas sur 15 pour la variante BatchNorm, 15 sur 25 pour le transformeur).

\subsubsection*{Deux conclusions du lot précédent ne survivent pas au changement de modèle}

Le critère de sortie du lot exigeait que toute conclusion « par type de couche » tienne sur au moins
deux modèles. Deux n'y parviennent pas, et sont consignées comme telles :

\begin{center}\small
\begin{tabular}{@{}>{\raggedright\arraybackslash}p{5.4cm}>{\raggedright\arraybackslash}p{2.0cm}>{\raggedright\arraybackslash}p{2.6cm}>{\raggedright\arraybackslash}p{2.6cm}@{}}
\toprule
Conclusion & A1 & A2 / A3 & Verdict \\
\midrule
Le raccourci diagonal d'AdaFisher coûte quelque chose (HF3) & confirmée & confirmée & \textbf{tient} \\
Le terme croisé $\gamma$--$\beta$ n'est pas négligeable (HF4) & confirmée & confirmée & \textbf{tient} \\
Toutes les approximations se dégradent au fil de l'entraînement & 95\,\% des cas & 52 à 71\,\% & \textbf{tendance, pas règle} \\
Les couches sont fortement couplées entre elles (HF5) & minimum $0{,}63$ & $0{,}28$ (transformeur), $0{,}44$ (convolutif) & \textbf{ne tient pas} \\
\bottomrule
\end{tabular}
\end{center}

La dernière ligne est instructive dans sa direction : les couches sont \emph{moins} liées entre
elles là où les poids sont partagés que dans un simple perceptron multicouche. Ignorer les
corrélations inter-couches — ce que font toutes les méthodes de ce dépôt — est donc \emph{mieux}
justifié sur une convolution ou un transformeur que sur le modèle jouet où on l'avait d'abord
mesuré : exactement l'inverse de ce qu'une étude limitée au perceptron aurait laissé croire.

Quant aux amplitudes, elles valent la peine d'être citées : le coût du raccourci diagonal atteint
$0{,}86$--$0{,}90$ sur les convolutions et $0{,}83$--$0{,}97$ sur les couches du transformeur en fin
d'entraînement (contre $0{,}25$--$0{,}41$ sur les têtes de classification), soit un ordre de grandeur
au-delà du seuil de 10\,\% fixé par l'hypothèse ; et le terme croisé des normalisations porte
$0{,}65$--$0{,}70$ de l'énergie du bloc sur le convolutif, $0{,}39$--$0{,}62$ sur le transformeur.

\subsection{Un problème numérique corrigé en cours de route : les valeurs propres exactement nulles}

Une conséquence directe et concrète du fait que $F$ possède des directions de rang nul (comme le
noyau de $d_{\rm in}+1$ dimensions de la tête de classification, ou les pixels de bordure
identiquement nuls de MNIST évoqués plus haut) s'est manifestée pendant les campagnes
d'entraînement elles-mêmes, avant même le début de cette campagne de mesure : la décomposition en
valeurs propres nécessaire aux modes \texttt{ekfac}/\texttt{tekfac} peut échouer brutalement sur le
GPU quand la matrice à décomposer est \emph{exactement} déficiente en rang (et non pas seulement
« mal conditionnée » au sens usuel), un cas que deux campagnes d'entraînement réelles ont
effectivement rencontré. La correction (un petit terme de régularisation ajouté avant la
décomposition, mathématiquement inerte sur le résultat final puisqu'il ne fait que décaler
uniformément les valeurs propres sans toucher aux vecteurs propres) est maintenant en place, et ce
même phénomène de rang exactement nul est aujourd'hui \emph{quantifié} plutôt que seulement contourné,
grâce au calcul dense complet décrit ci-dessus.

% =====================================================================
\section{Feuille de route : les lots d'implémentation restants}
% =====================================================================

Le travail de cette campagne est découpé en lots successifs, chacun construit sur les précédents et
validé par ses propres tests avant de passer au suivant — le même principe de « une itération à la
fois » qui a déjà structuré la construction des cinq modes de l'optimiseur.

\begin{center}
\small
\begin{tabular}{@{}p{0.6cm}p{5.4cm}p{5.6cm}@{}}
\toprule
Lot & Contenu & Où ça en est \\
\midrule
0 & Les fondations : politique de précision numérique, désactivation du mode de calcul approché
\eng{TF32}, ensembles d'exemples sondes reproductibles, registre qui classe chaque paramètre de
chaque modèle par type de couche, et le pont qui sait lire les instantanés déjà sauvegardés par les
bancs d'essai. & \textbf{Terminé} : 27 tests dédiés, aucun changement au code existant. \\
1 & Le premier modèle traité de bout en bout, en régime dense : capture des gradients par exemple,
les trois sources (type-2, Monte-Carlo, empirique), et le calcul exact de $F$, $\hat E$ et des blocs
par couche, en double précision. & \textbf{Terminé, sauf le calcul complet à grande échelle}, qui
est lui-même désormais terminé (section~\ref{sec:mesures}). \\
2 & Le « zoo » complet des approximations (bloc-diagonal, K-FAC, EKFAC, TKFAC, diagonale) et les
métriques M1, M3, M5, M7, M8, appliqués au premier modèle, sur ses cinq instantanés ; premier tracé
du plancher de bruit. & \textbf{Terminé} (section~\ref{sec:lot2}) ; deux lectures de ses résultats
ont dû être corrigées par le lot 3. \\
3 & Extension à un modèle avec convolutions et normalisations (dans ses deux variantes), et à un
modèle avec attention ; la décomposition partage de poids / indépendance ; intervalles de confiance
par couche. & \textbf{Terminé} (section~\ref{sec:lot3}) : HF2 tranchée sur trois modèles, trois
erreurs de code antérieures corrigées au passage, et deux conclusions du lot 2 invalidées. \\
4 & Le régime B (couche par couche, sans jamais former la grande matrice), validé d'abord contre le
régime dense, puis appliqué aux modèles plus grands — dont \textbf{l'auto-encodeur MNIST et son
blocage} (section~\ref{sec:stall}). & À venir ; c'est le lot qui doit trancher la cause du blocage. \\
5 & Le protocole opérationnel P2 : lecture directe de l'état réel de l'optimiseur, avec le
réchauffage déjà quantifié (section~\ref{sec:protocoles}). & À venir. \\
6 & Extensions optionnelles : un modèle à largeur réaliste, un modèle en régime sans matrice, et
d'éventuels nouveaux modèles pour couvrir les types de couches encore absents (\eng{embeddings} de
texte, LoRA). & Optionnel, décidé après le lot 5. \\
\bottomrule
\end{tabular}
\end{center}

L'ordre n'est pas arbitraire : le régime dense (lot 1--3) sert de \textbf{juge de référence} pour
valider le régime factorisé du lot 4 avant de lui faire confiance sur des modèles où aucun calcul
dense n'est possible ; et le protocole opérationnel (lot 5) a besoin du zoo d'approximations déjà
construit et validé par le lot 2. Créer un tout nouveau modèle de banc d'essai (lot 6) est
délibérément placé en dernier, car cela représenterait un changement dans les bancs d'essai
eux-mêmes, que cette campagne s'est fixé comme règle de ne faire qu'en tout dernier recours.

% =====================================================================
\section{Risques identifiés et façons de les traiter}
% =====================================================================

\begin{itemize}
\item \textbf{Le rang déficient du régime B.} Comme expliqué section~\ref{sec:regimes}, dans le
régime factorisé, toute métrique qui passe par un inverse mélange une vraie erreur d'approximation
avec de la courbure « inventée » là où la référence n'a simplement aucune information. Parade :
toujours présenter ces métriques comme des courbes en fonction de l'amortissement $\lambda$, jamais
comme un point unique, et confronter au régime dense chaque fois que c'est possible.
\item \textbf{La précision numérique.} Les spectres de courbure s'étalent sur plus de huit ordres de
grandeur ; toutes les références et les Grams intermédiaires sont donc calculés en double précision
(\eng{fp64}), et le mode de calcul approché \eng{TF32} des GPU récents (qui introduirait une erreur
relative de l'ordre de $10^{-3}$, suffisante pour fausser tous les tests d'exactitude) est
explicitement désactivé et vérifié par un test dédié.
\item \textbf{Un piège matériel concret.} Les tâches d'entraînement du dépôt tournent sur une petite
tranche de GPU (10 Go de mémoire) ; une tâche d'analyse qui hériterait de cette même contrainte de
mémoire ne pourrait tout simplement pas former $F$ pour les modèles en régime dense. Chaque tâche
d'analyse générée déclare donc explicitement son propre besoin en mémoire, plus généreux.
\item \textbf{Des instantanés mal étiquetés.} Une correction apportée en cours de campagne au calcul
des fractions d'entraînement (0\,\%, 1\,\%, \dots) a rendu incorrects tous les instantanés produits
\emph{avant} cette correction, à l'exception de ceux du mode \texttt{diag}. Le pont de lecture des
instantanés refuse activement de charger un instantané dont les métadonnées indiquent qu'il provient
d'avant cette correction, plutôt que de risquer une comparaison silencieusement faussée.
\item \textbf{Des graines encore en cours d'exécution.} Comme signalé section~\ref{sec:checkpoints},
l'analyse doit systématiquement rapporter le nombre de graines réellement disponibles au moment de
la mesure, jamais un nombre visé par avance.
\item \textbf{Le volume de stockage.} Les matrices intermédiaires du régime B (les « Grams » par
couche) ne sont jamais écrites sur disque, seules les métriques finales le sont — sans quoi le
volume de données produites deviendrait rapidement ingérable.
\end{itemize}

% =====================================================================
\section{Ce que la campagne produira, et à quel coût}
% =====================================================================

\subsection{Le format des résultats}

Chaque combinaison (modèle, bras, graine, fraction d'entraînement) produit un tableau de résultats
au format \eng{CSV}, avec une ligne par mesure (une couche, une structure d'approximation, une
source, une référence, une valeur d'amortissement, une métrique, sa valeur) et un petit fichier
d'accompagnement qui trace, pour chaque tableau, la version exacte du code utilisée, les exemples
sondes précis employés, et les métadonnées de l'instantané d'origine — de sorte que tout résultat
reste entièrement reproductible et vérifiable après coup.

\subsection{Les figures prévues}

\begin{itemize}
\item une carte de l'efficacité $\rho$, croisant chaque structure d'approximation et chaque type de
couche, une par référence ($F$ ou $\hat E$) ;
\item un diagramme en barres décomposant l'erreur totale en ses trois composantes (source,
structure, interaction) ;
\item des courbes suivant l'évolution de la fidélité au fil de l'entraînement, séparément pour les
sondes d'entraînement et de validation ;
\item une carte du biais d'indépendance ($\sigma_2/\sigma_1$) par couche ;
\item des matrices de couplage entre couches différentes ;
\item un nuage de points comparant, couche par couche, le classement obtenu en conditions idéalisées
(P1) contre le classement obtenu en conditions réelles (P2) ;
\item une courbe du plancher de bruit en fonction du nombre d'exemples sondes utilisés ;
\item un panneau dédié spécifiquement aux couches de normalisation.
\end{itemize}

\subsection{Le budget de calcul}

Le calcul dense sur le premier modèle (A1) a été mesuré à un peu plus d'une heure sur une tranche de
GPU, avec un pic de mémoire de 6,06 Go — très proche de la prévision théorique de 6,0 Go faite avant
de lancer la tâche, ce qui donne confiance dans les estimations de coût pour les modèles restants.
La campagne principale complète (les modèles en régime dense et factorisé, cinq instantanés, les
graines disponibles, les deux protocoles P1 et P2) est estimée à l'ordre de 100 à 200 heures de
calcul GPU sur la tranche de matériel utilisée pour l'analyse.

% =====================================================================
\section{Glossaire rapide}
% =====================================================================

\begin{longtable}{@{}p{3.4cm}p{11.6cm}@{}}
\toprule
Terme & Signification en clair \\
\midrule
\endhead
Matrice de Fisher $F$ & la courbure exacte du problème d'apprentissage, moyennée sur les étiquettes
que le modèle imagine lui-même possibles. \\
Fisher empirique $\hat E$ & la même idée, mais moyennée seulement sur les étiquettes réellement
observées dans les données — bien moins chère à calculer, structurellement biaisée. \\
GGN (Gauss--Newton généralisée) & une autre courbure classique, mathématiquement identique à $F$
dans tous les cas pratiques utilisés dans ce dépôt. \\
K-FAC & une façon d'approcher la courbure d'une couche en la factorisant en un produit de deux
matrices bien plus petites (une pour les entrées de la couche, une pour ses gradients de sortie). \\
EKFAC & K-FAC, en réajustant après coup l'importance de chaque direction de façon optimale. \\
TKFAC / TEKFAC & une autre factorisation, construite pour préserver exactement l'énergie totale
(la trace) de la courbure exacte. \\
Amortissement ($\lambda$) & un petit nombre ajouté à la diagonale d'une matrice pour pouvoir
l'inverser sans risque, même si certaines de ses valeurs propres sont nulles ou très petites. \\
EMA (moyenne mobile exponentielle) & une façon de lisser une mesure dans le temps, en mélangeant
la nouvelle observation avec la mémoire des observations précédentes. \\
Régime A / B / C & trois façons de calculer la courbure exacte selon la taille du modèle : tout
former explicitement (A), ne jamais former la grande matrice mais garder des Grams par couche (B),
ou ne calculer que des produits matrice-vecteur sans jamais rien former (C). \\
Protocole P1 & mesure de la meilleure fidélité \emph{possible} d'une approximation, dans des
conditions idéalisées. \\
Protocole P2 & mesure de la fidélité \emph{réellement obtenue} pendant un entraînement réel, avec
tous les compromis pratiques empilés. \\
Réchauffage (\eng{re-warm}) & le fait de relancer un optimiseur neuf pendant quelques centaines de
pas, à partir d'un instantané qui ne contient que les poids, pour reconstituer une mémoire interne
représentative avant de la mesurer. \\
Plancher de bruit & la variabilité statistique minimale, inévitable, d'une mesure faite sur un
nombre fini d'exemples — en dessous de laquelle une différence n'est pas significative. \\
\eng{Checkpoint} / instantané & une sauvegarde des poids d'un réseau à un instant donné de son
entraînement. \\
\eng{Arm} / bras & une des sept stratégies d'entraînement comparées (les cinq modes de courbure,
plus Adam, plus AdamW). \\
\bottomrule
\end{longtable}

\end{document}
